\section{The assembled evidence for the coverage model}\label{app:evidence}

Table~\ref{tab:evidence} collects the evidence for the coverage model in one place, with each claim's
standing stated explicitly, since the support is deliberately heterogeneous: a representation
theorem, certified measurements, a confirmed prediction, and a set of benchmark results. Detailed
treatments are at the listed locations.

\begin{table}[h]
  \centering
  \footnotesize
  \begin{tabular}{@{}p{0.26\linewidth}p{0.30\linewidth}p{0.18\linewidth}p{0.17\linewidth}@{}}
    \toprule
    Claim & Evidence & Standing & Where \\
    \midrule
    The ceiling $c$ exists & bounded loss with decaying marginal damage on exact lattices & measured & Section~\ref{sec:coverage} \\
    The class $g\bigl(\sum_i b_i\bigr)$ is canonical & Acz\'el--Ling representation & proven & App.~\ref{app:link} \\
    $f$ is within ${\sim}2\%$ ordinal defect of the class & split-half certification against matched-noise nulls & measured, certified & Section~\ref{sec:ordinal}, App.~\ref{app:ordinal} \\
    Alternating M\"obius spectrum, overshoot up to $32\times$ & \thmthree{} prediction confirmed on exact lattices & confirmed prediction & Section~\ref{sec:duality}, App.~\ref{app:proofs} \\
    Certificate identity and its $O(L)$ product form & \thmone{}, \thmfour{}, \propfive{} & proven, numerically verified & Sections~\ref{sec:walsh}, \ref{sec:tau}, App.~\ref{app:proofs} \\
    $\tau$ organizes the measured order-$\ge 2$ shares & 36 dense-model and 32 MoE-block cells; Spearman $0.958$ ($\tau$), Pearson $0.984$--$0.997$ (exact form) & measured & Section~\ref{sec:tauval} \\
    The share is forecastable from sampled configurations & estimator comparison against held-out lattices (Table~\ref{tab:certest}) & measured & Section~\ref{sec:tauval} \\
    The generator is probit-$\sqrt{\cdot}$, with exp its tractable approximation & slope calibration across six cells & measured & App.~\ref{app:knapsack} \\
    Aumann--Shapley recovers the optimal coefficient at an effective density, and its layer ordering exactly & Prop.~B.1; real-model correlation $0.785$ & proven under coverage; empirical validation & App.~\ref{app:as} \\
    The fit transports across deployment densities & degradation $1.3$--$2.0\times$ against the GP's $2.1$--$4.9\times$ at $N=L/2$ & measured & Section~\ref{sec:predictor} \\
    \bottomrule
  \end{tabular}
  \caption{The evidence for the coverage model, with each claim's standing. Proven means a theorem
  (numerically verified where marked),
  measured means a lattice or benchmark result, and a confirmed prediction is a property of the data
  that the fit did not use.}
  \label{tab:evidence}
\end{table}

\section{The model class $g\bigl(\sum_i b_i\bigr)$: canonicity, and details for the ordinal
analysis}\label{app:link}

The Walsh--Hadamard expansion of $f$ is exactly the functional-ANOVA decomposition on the
$p$-biased Boolean cube: order 0 is the mean, order 1 the per-layer main effects, order $k$ the $k$-way
components. The tempting reading, that order-$\ge 2$ energy measures interaction among layers, does
not follow, because ANOVA additivity is additivity in the loss values and the loss saturates, so a model
that is additive in the loss must spend order-$\ge 2$ energy merely to represent the curvature of
saturation, even with no coupling among layers. Order-$\ge 2$ energy therefore conflates real coupling
with the shape of a scalar transformation. The class $L(S)=g\bigl(\sum_{i\in S} b_i\bigr)$, additive after
a strictly increasing $g$, partially disentangles the two, since if a single such $g$ linearizes $f$ (that is,
$g^{-1}\circ f$ is order-1), the order-$\ge 2$ energy was produced by $g$ rather than by interaction among
layers.
% Coverage is exactly such a model, additive in $b$ with a saturating generator, and \thmfour{}
% together with the collapse of Section~\ref{sec:tauval} shows that the loss-space order-$\ge 2$ energy is
% predicted by the generator. The same accounting explains the collapse below 4 bits described in
% Section~\ref{sec:allocsub4}, where
% the generator is near-vertical, so order-$\ge 2$ energy explodes (held-out additive $R^2=-0.92$) without
% new structure among layers appearing. (The GPTQ ladder of Section~\ref{sec:allocsub4} behaves
% differently: the coverage-form certificate forecasts at most 0.02 there, yet the additive prediction
% underestimates the allocator's chosen configurations by 16--57\%.) 
The residual test that survives this critique is the order-$\ge 2$
energy after the best transformation, measured in Section~\ref{sec:ordinal}, with protocol details below.

\paragraph{Why searching over transformations is canonical.} The class $g\bigl(\sum_i b_i\bigr)$ is not a
parametric convenience. It is exactly the class of damage functions that accumulate one layer at a time by
a fixed rule, $f(S\cup\{i\}) = f(S) \oplus d_i$, with a per-layer increment $d_i$ that does not depend on
$S$ and an operation $\oplus$ that is associative, continuous, and strictly increasing with identity $0$.
By the Acz\'el--Ling
representation theorem, every such $\oplus$ is conjugate to ordinary addition: there is a strictly
increasing generator $g$ with $x\oplus y = g\bigl(g^{-1}(x)+g^{-1}(y)\bigr)$, unique up to scale
\citep{aczel1966fe,ling1965assoc}. Accumulating damage by any rule with these properties is therefore
equivalent to $f$ being additive after some strictly increasing transformation, and the search of
App.~\ref{app:knapsack} ranges over such
rules. Coverage is the instance $x\oplus y = x+y-xy/c$, the probabilistic sum (noisy-OR)
\citep{pearl1988probabilistic,klement2000tnorms}, whose generator is the log-headroom
$g^{-1}(f) = -\log(1-f/c)$. 
The essential interaction of Section~\ref{sec:ordinal} is then precisely what no rule with
these properties can absorb,
namely the ordinal defect.
% The probit-$\sqrt{\cdot}$ rule of App.~\ref{app:knapsack} is
% another instance. 

\subsection{Details for Section~\ref{sec:walsh}: the ordinal certification and the noise-corrected
energies}\label{app:ordinal}

\begin{proof}[Proof of \thmtwo]
($\Rightarrow$) A strictly increasing $g$ preserves and reflects strict order. ($\Leftarrow$) On the finite
set $V=\{b\cdot x\}$ define $g_0(b\cdot x):=f(x)$. This is well-defined (if $b\cdot x=b\cdot y$, the order
condition forces $f(x)=f(y)$) and strictly increasing (if $b\cdot x<b\cdot y$ then $f(x)<f(y)$), and it
extends piecewise-linearly to a strictly increasing $g$ on $\mathbb{R}$.
\end{proof}

The claim is the finite binary case of additive conjoint measurement
\citep{lucetukey1964conjoint,debreu1960topological,scott1964measurement}, and is stated for
completeness.

\paragraph{Certification protocol: the estimand and the noise floor.} This and the next three
paragraphs detail
how the ordinal defect of Section~\ref{sec:ordinal} is measured and certified. The estimand is the
draw-expected damage $\tilde f$: where the measured $f(S)$ of Section~\ref{sec:walsh} averages 16
stochastic-rounding draws, $\tilde f(S)$ is its expectation over draws, and all randomness below is over
those draws. The ranking constraints
of $\tilde f$ are the inequalities $b\cdot(x-y)>0$, one per ordered pair of configurations with
$\tilde f(x)>\tilde f(y)$. By
\thmtwo{}, an index satisfying every constraint exists exactly when $\tilde f$ is additive up to a strictly
increasing transformation, and the ordinal defect is
the smallest fraction of the constraints that any index must violate. The complication is that $\tilde f$
is known only up to the draw noise. Both
bounds below therefore share two devices. First, the draws are split into two halves, whose means
$\bar f_A$ and $\bar f_B$ are independent estimates of $\tilde f$ at every configuration, the split-half device
of reliability theory \citep{spearman1910correlation,brown1910some}. Second, a pair of
configurations contributes a constraint on a given half (is \emph{eligible} there) only if its gap on
that half exceeds twice its standard error. A half-mean at configuration $S$ has standard error
$\hat\sigma(S)/\sqrt{8}$, with $\hat\sigma(S)$ the sample standard deviation of the draws, so the gap
between configurations $x$ and $y$ has standard error
$\sqrt{\hat\sigma(x)^2 + \hat\sigma(y)^2}/\sqrt{8}$.

\paragraph{The upper bound: an exhibited index.} Because the defect is a minimum over indices, the
fraction of constraints violated by any single index bounds it from above. The tightest such bound would
come from the index minimizing the violated fraction, which is NP-hard to find \citep{amaldi1995maxfs},
so we exhibit instead the solution $(b, \xi)$ of the hinge relaxation
\[
\min_{b,\,\xi} \ \sum_p \xi_p \qquad \text{subject to} \qquad
b\cdot(x_p - y_p) \ \ge\ 1 - \xi_p, \quad \xi_p \ge 0 ,
\]
the pairwise-preference construction of ranking with slack variables \citep{joachims2002ranking}, here
without the norm term, so the problem is a linear program, fit on the eligible constraints of half A. The
fitted index's violated fraction is then evaluated twice: in-sample, on the half-A constraints it was fit
on, and cross-half, on half B's own eligible constraints, which the fit never saw. The two evaluations
play different roles. 
% If $f$ were measured exactly, no split would be needed, since the in-sample
% fraction of a fit on all constraints would be both a valid upper bound and the tightest one this index
% family offers. 
With noisy draws the reported number is instead an estimate of the index's violation rate
on $\tilde f$'s constraints, and the in-sample rate is biased low for that quantity. The
cross-half rate removes this adaptation, evaluating a fixed index on independently selected constraints.
We report both and quote the larger. On most cells the two agree to within about a tenth of a percentage
point.

\paragraph{The lower bound: certified witnesses.} The witnesses are rank flips, two contexts that order
the same two layers oppositely (the infeasibility witnesses of Section~\ref{sec:ordinal}). For a layer pair
$(i,j)$ and a context $S$ containing neither, the constraint between $\tilde f(S\cup\{i\})$ and
$\tilde f(S\cup\{j\})$
involves $b$ only through the sign of $b_i - b_j$, which is the same for every context. If, among
$\tilde f$'s constraints, $m$ contexts order $i$ above $j$ and $n$ order $j$ above $i$, then every additive index
violates at least $\min(m,n)$ of them. Summing $\min(m,n)$ over layer pairs, whose constraint families are disjoint,
counts constraints that every index must violate, so this count divided by the number of eligible
constraints is a lower bound on the defect, and by \thmtwo{} a
single opposed pair of contexts already refutes every increasing transformation. To certify a context
against noise, it must be eligible on half A and recur with the same sign on half B. The witness
count
of Table~\ref{tab:ordinal} is the sum over layer pairs of the smaller number of certified contexts of
each preference direction.

\paragraph{The matched-noise null.} The noise floor and the eligible constraint check still pass a small number of
false positives, so every statistic above, for both bounds, is read against a matched-noise null lattice
pushed through the identical pipeline. The null is the coverage model fitted to the measured means, which
is additive after an increasing transformation and so contains no true flips, with the 16 draws
regenerated as independent Gaussian samples centered on the fit, scaled by each configuration's measured
per-draw sample standard deviation $\hat\sigma(S)$. This is a parametric bootstrap under the hypothesis of no
essential interaction \citep[Section~6.5]{efron1993bootstrap}. The Gaussian form is an assumption, and a mild
one here, because every statistic consumes the draws only through the two half-means, each an average of
8 draws that are themselves means of per-token losses.
%  so the half-means are close to Gaussian whatever
% the shape of a single draw. 
The null measures the false-positive rate of the whole procedure, and a cell
is certified only where the measured statistic exceeds its null.

\begin{table}[t]
    \centering
    \small
    \begin{tabular}{llccccccc}
      \toprule
      & & \multicolumn{2}{c}{Witnesses} & \multicolumn{2}{c}{Sign replication} & \multicolumn{2}{c}{LP upper bound (\%)} & $E^*$ (\%) \\
      \cmidrule(lr){3-4}\cmidrule(lr){5-6}\cmidrule(lr){7-8}\cmidrule(lr){9-9}
      Model & Activation scales & meas. & null & meas. & null & meas. & null & \\
      \midrule
      Qwen3-0.6B & per-block & 13 & 18 & 0.76 & 0.66 & 0.6 & 0.8--1.1 & 4.9 \\
      Qwen3-8B (fp32) & per-block & 5 & 10 & 0.67 & 0.61 & 0.5--0.6 & 0.7--1.1 & 0.2 \\
      Llama-3.2-3B & per-block & 56 & 6 & 0.80 & 0.69 & 1.2--1.4 & 0.3--0.4 & 1.7 \\
      Qwen3-30B-A3B & per-block & 21 & 7 & 0.79 & 0.53 & 0.5--0.8 & 0.4--0.8 & $-0.4$ \\
      \midrule
      Qwen3-0.6B & per-tensor & 74 & 18 & 0.95 & 0.77 & 1.6 & 2.5--2.8 & 3.5 \\
      \bottomrule
    \end{tabular}
    \caption{Certification of the ordinal defect on full $2^8$ lattices. 
    Only comparisons whose loss gap exceeds twice its standard error enter any
    count. Witnesses count is the sum over layer-pairs of the smaller number of contexts of each preference direction
    that sign-replicate across the two halves of the draws. Sign replication is the replicated fraction of half-A comparisons.
    LP upper bound is the fraction of eligible ranking constraints
    violated by the additive index fit on half A, evaluated in-sample on half A and cross-half on half B.
    The range spans the two evaluations, and the text quotes the larger. Every statistic is compared
    with the identical statistic on the matched-noise null, which contains no true flips. $E^*$ is the
    order-$\ge 2$ energy share remaining after the best strictly increasing transformation of the loss
    (the free-monotone estimator of App.~\ref{app:ordinal}).}
    \label{tab:ordinal}
  \end{table}

\paragraph{Per-model detail.} Table~\ref{tab:ordinal} reports every cell against its null. In the deployed
per-block regime, Qwen3-0.6B and Qwen3-8B sit at or below their nulls and are
statistically indistinguishable from zero, Qwen3-30B-A3B has a defect of at most $0.8\%$, at its null, and
Llama-3.2-3B carries a small genuine residue (56 witnesses against a null of 6). In the per-tensor regime
the witnesses exceed their null (74 against 18, at $0.95$ sign replication), which is the outlier
coupling of Section~\ref{sec:scope} appearing as certified rank reversals. An additive index nevertheless still
orders at least $98\%$ of the ranking constraints above the noise floor correctly. Nearly all of the apparent high-order energy in this
worst cell is therefore curvature of the loss coordinates, as seen by $E^\star$.

\paragraph{Removing the stochastic-rounding noise floor from measured order energies.} The split of the
draws serves a second, separate role for the order energies of Section~\ref{sec:walsh}, where the halves are
differenced rather than held out. Each measured
$f(S)$ averages $K$ stochastic-rounding draws and so equals the draw-expected damage $\tilde f(S)$
plus a mean-zero
Monte-Carlo error with variance $\sigma^2(S)/K$, where $\sigma^2(S)$ is the per-draw variance at
configuration $S$. The Walsh transform is linear, so the measured coefficients are likewise unbiased.
Their squares are not, since
$\E[(a+Z)^2] = a^2 + \Var(Z)$ for constant $a$ and mean-zero $Z$, so each measured order energy $W_k$
exceeds the energy of the draw-expectation by the summed error variances of its order-$k$ coefficients,
a floor
$\nu_k > 0$. The per-coefficient error variance is nearly the same for every basis function
(independent per-configuration noise is white in an orthonormal basis: exactly even at $p=\tfrac12$,
within a small per-order factor at $p=0.6$), so $\nu_k$ is close to proportional to the number
$\binom{L}{k}$ of order-$k$ basis functions. All but $L$ of the $2^L-1$ non-constant functions have
order $\ge 2$, so the floor's mass lands almost entirely in the order-$\ge 2$ energies and deflates the
measured order-1 share. To estimate it, we average two disjoint halves of the draws
separately, $\bar f_A$ and $\bar f_B$, and form the difference table
$d(S)=\tfrac12\bigl(\bar f_A(S)-\bar f_B(S)\bigr)$, in which the draw-expectation cancels exactly and a
pure mean-zero error of the same variance $\sigma^2(S)/K$ remains, so the order energies of $d$ satisfy
$\E[W^{(d)}_k]=\nu_k$. This requires only that the two halves are independent (fresh rounding draws per
configuration and draw). The corrected
energies reported in Sections~\ref{sec:walsh} and \ref{sec:tauval} are
$\max(W_k - W^{(d)}_k,\,0)$. Subtracting replicate-estimated error variance from an observed second
moment is the classical variance-components device of the analysis of variance
\citep[Ch.~7]{scheffe1959anova}, applied here to split-half replicates
\citep{spearman1910correlation,brown1910some}.

\paragraph{Estimating the transformation-independent energy floor.} For strictly increasing $g$ let $E(g)$
be the relative order-$\ge 2$ Walsh energy of the transformed damage $g^{-1}\circ \tilde f$. By \thmone{}
applied
to $g^{-1}\circ \tilde f$, $E(g)$ is the fraction of the transformed damage's variance that the best additive
model leaves unexplained, and $E^* := \inf_g E(g)$ is the floor of that fraction over all strictly
increasing transformations. A small $E^*$
therefore certifies that $g^{-1}\circ \tilde f$ is nearly additive for some strictly increasing $g$, so the
measured high-order energy is curvature of the loss coordinates rather than coupling among layers.
Estimation must respect two hazards: Monte-Carlo noise spreads nearly evenly across the Walsh basis, as
above (247 of 255 basis functions have order $\ge 2$, so raw noise masquerades as high-order structure),
and $g^{-1}$ is nonlinear, amplifying noise where the inverse is steep. The protocol therefore
splits the 16 rounding draws into three parts and fits every transform on part 1 only. The parametric
families are fit by nonlinear least squares, and the free monotone transform by alternating a weighted
degree-$\le 1$ fit with isotonic regression against the ordering of $f$'s values, which searches all
increasing transformations at once and is the fit that estimates $E^*$: isotonic regression returns the
least-squares best monotone nondecreasing fit to a scatter, computed exactly by the
pool-adjacent-violators algorithm \citep{ayer1955pava,barlow1972isotonic}. Each squared Walsh coefficient of
the transformed damage is then estimated by multiplying the coefficient computed from part 2 by the
coefficient computed from part 3, an unbiased estimate of the squared coefficient because the two parts'
noises are independent and mean zero. The product is unbiased but not constrained to be nonnegative,
so a cell whose true share is near zero can estimate slightly negative, as the $-0.4\%$ entry of
Table~\ref{tab:ordinal} does. The headline numbers appear in Section~\ref{sec:ordinal}, and the
per-cell values in the last column of Table~\ref{tab:ordinal}. 

% The estimate
% under a fitted parametric family is fragile, because a misfitted ceiling makes the inverse near-vertical
% at the top of the range and manufactures high-order energy ($4.9\%\to 28\%$ on the 0.6B cell under a
% misfit exponential ceiling), so the quoted $E^*$ comes from the free-monotone transform.

\paragraph{What depends on the measure.} Being exactly degree-$\le 1$, raw or after a transformation, is
an algebraic property of $f$ (of its unique multilinear extension), independent of the product measure,
since the span of the degree-$\le k$ $\mup$-Walsh functions is the same subspace for every $p$. What depends
on $p$ is the accounting: the best degree-$\le 1$ projection, its coefficients $w_i(p)$, and the
order-$\ge 2$ energy are all priced under the chosen measure. Two consequences follow. \thmone's
certificate is stated, and must be read, at the deployment $\mup$. The wrong anchor of Section~\ref{sec:anchor} is
ultimately a wrong measure, the isolated slope $s_i$ being the $p\to 0$ limit of the deployment-point slope
$w_i(p)$. 
% And the coverage fit transports across $p$ (Section~\ref{sec:predictor}, Fig.~\ref{fig:ood}) while
% loss-scale additive coefficients must shift with the operating point.

\section{Experimental setup}\label{app:setup}

We collect the protocol shared across Sections~\ref{sec:walsh}--\ref{sec:scope} here. Section-specific
deviations are noted in place.

\paragraph{Models.} \emph{LLMs}: Qwen3-0.6B, Qwen2.5-3B, Qwen3-4B, Qwen3-8B, Llama-3.2-3B,
 with the per-linear order-share cells of
Section~\ref{sec:walsh} extending to Qwen3-14B, Qwen3-32B, Llama-3.3-70B, Qwen2.5-72B, and
gpt-oss-120B; Qwen3-30B-A3B, Qwen3-235B-A22B, and GLM-4.6 (355B
MoE) for allocation at scale; Qwen3.6-27B for the vendor-checkpoint serving comparison,
with damage measured and ModelOpt calibrated on a chat-templated task corpus (96 MATH-train, 64
gsm8k-train, 96 MMLU-Pro validation prompts); Qwen3-235B and GLM-5.2 (753B) for the
scale-robustness checks.
\emph{Diffusion}: Qwen-Image-2512 DiT ($L=60$ transformer blocks) for the diffusion
prediction results. All weights are the public checkpoints.

\paragraph{Quantization formats.} The measurement regime throughout is W4A4 with NVFP4 per-block
scales. We also analyze W16A4 (activation-only, isolating the channel that carries the
non-additivity) and per-tensor activation scaling (a coarse setting that induces interactions of every
order among layers). The allocation alphabet is $\{16, 8, 4\}$ bits (bf16 / FP8 / NVFP4) for the
servable studies and is extended to $\{16,8,4,2,1\}$ for the
below-4-bit collapse study. Serving uses three pipelines. NVIDIA ModelOpt PTQ with TensorRT-LLM
serves the $\ge 4$-bit MoE deployment recipes, and vLLM the vendor-checkpoint comparison of
Section~\ref{sec:alloc27b}; these paths serve W4A4 ModelOpt exports with runtime activation
quantization on the NVFP4 units (NVIDIA's published 27B checkpoint, which keeps its NVFP4 linears
weight-only, is the exception), as does the SGLang-served GLM-5.2 comparison. 
% The served-accuracy
% grids (Tables~\ref{tab:ge4served30},
% \ref{tab:ge4served235}, \ref{tab:sub4served30}, and \ref{tab:sub4code}) instead realize the weight
% channel only: weights are cast to each configuration's assigned formats (round-to-nearest at
% $\ge 4$ bits, GPTQ below four), exported as ordinary bf16 checkpoints, and served by SGLang with
% activations left unquantized, so these tables test the allocation's weight placement rather than
% the full W4A4 pipeline. A W4A4 replication of the 30B grid subset (Section~\ref{sec:allocmoe}:
% additive, coverage, and ModelOpt at effective bits 4.5 and 5.0 as ModelOpt exports with runtime
% activation quantization, TensorRT-LLM, AIME-2025 over 8 seeds) gives 71.7/69.6/70.8 at eff 4.5 and
% 69.2/72.9/68.8 at 5.0, against the weight-only grid's 75.4/69.6/70.0 and 69.2/69.6/72.9.
\paragraph{The damage set function and the lattice.} For a contiguous block of $L$ layers we form
$f(S)=\lce(S\ \text{quantized}) - \lce(\text{bf16})$ by running the model with exactly the units in $S$
quantized. Every exact lattice in the paper measures such a block centered at the network's mid-depth (block
offset $\lfloor(L'-L)/2\rfloor$ in an $L'$-layer network, so the middle eight of Qwen3-0.6B's 28
layers), except the Qwen3-30B-A3B extension of Fig.~\ref{fig:tau}, which repeats the measurement at
eight block positions. Restricting to a small block is what makes the $2^L$ enumeration feasible, and
the enumerated lattice is a reference against which sampled estimates, fitted models, and the
noise corrections of App.~\ref{app:ordinal} are certified. The full-network cells
(Section~\ref{sec:walsh}, the prediction benchmark, and the allocation studies) instead quantize units
anywhere in the network, where enumeration is impossible, and sample configurations from the $p$-biased
measure $\mup$. Each $f(S)$ is
averaged over stochastic-rounding draws (3 seeds for the headline sampled cells, std $\le 0.03$; 8--16
draws for the exact lattices of Section~\ref{sec:walsh}, whose noise treatment is in App.~\ref{app:ordinal}).
The
Walsh/M\"obius transforms are computed from measured $f(S)$. The order-$\ge 2$
energy certificate is the measured $\sum_{k\ge 2}W_k/\Var(f)$ where a lattice is affordable and, at scale,
the closed-form forecast from the fitted break-rates, namely the exact heterogeneous product form
$1 - \sum\beta_i/[\prod(1+\beta_i)-1]$ of \propfive{} (Section~\ref{sec:tau}), whose symmetric summary is $\tau$
(share $\approx\tfrac12\tau^2$). The ``$\tau$-forecast'' label refers to this pipeline. Where
a measured share exists for comparison, the forecast understates it (Qwen3-30B: forecast
$3{\times}10^{-4}$ against measured 0.013). Even at that ratio the 235B and GLM-4.6 forecasts
(${\sim}10^{-4}$) stay below 0.01, so they certify additive dominance, not a precise share.

\paragraph{Calibration and evaluation data.} Prediction and interaction-order studies use WikiText-2
(``wiki'') and a Python code corpus, CodeSearchNet functions (``code''). 
The MoE allocation studies use a task-aligned calibration set: BFCL-v3
(live\_simple, live\_multiple, simple\_python, multiple) tool-use prompts plus RULER validation prompts,
tokenized to 16 sequences of 2{,}048 tokens. The fake-quant metric is the KL divergence
$\mathrm{KL}(\text{bf16} \,\|\, \text{quantized})$ of next-token distributions on the disjoint eval half
of this set. Served downstream evaluation uses AIME-2025, GPQA
diamond, MMLU-Pro, HumanEval, and gsm8k (Section~\ref{sec:allocation}).

\begin{table}[t]
  \centering
  \small
  \begin{tabular}{llcccccccc}
    \toprule
    & & \multicolumn{8}{c}{Context size $|S|$ (layers already quantized)} \\
    \cmidrule(lr){3-10}
    Model & Activation scales & 0 & 1 & 2 & 3 & 4 & 5 & 6 & 7 \\
    \midrule
    Qwen3-0.6B & per-block & .183 & .022 & .024 & .022 & .021 & .023 & .022 & .012 \\
    Qwen3-8B (fp32) & per-block & .023 & .006 & .006 & .006 & .006 & .007 & .003 & .009 \\
    Llama-3.2-3B & per-block & .016 & .013 & .010 & .009 & .009 & .007 & .009 & .003 \\
    Qwen3-30B-A3B & per-block & .060 & .027 & .024 & .023 & .022 & .016 & .021 & .021 \\
    Qwen3-0.6B & per-tensor & 10.6 & .997 & .381 & .173 & .111 & .078 & .048 & .011 \\
    \bottomrule
  \end{tabular}
  \caption{Saturation measured directly: mean marginal damage by context size on the certification
  cells of Section~\ref{sec:ordinal}. Each
  entry averages $f(S\cup\{i\}) - f(S)$ over the layer-context pairs with $|S|$ layers already
  quantized. In every cell the marginal damage is largest at $|S|=0$ and falls once other layers are
  quantized concurrently.}
  \label{tab:marginals}
\end{table}

\paragraph{Predictor fitting.} The coverage model $(c, \{a_i\})$ and the additive model are fit by joint
least squares rather than per-coefficient sampling, for the variance and consistency reasons given in
the next two paragraphs. Fits use $N$
sampled configurations. Fits are initialized by linearizing through the generator
(ordinary least squares on the log-headroom values at a provisional ceiling) and refined by bounded
trust-region least squares in loss
space. Fitting in loss space avoids the alternative of transforming the data and then regressing, which
is biased because noise pushed through the nonlinear transform no longer averages to zero.
The surrogate baselines (RF/MiCo, GP-RBF/AutoQRA, RBF kernel
regression/AMQ, CoopQ pairwise) are fit on the same configurations, and GP/RF are
additionally converted into allocators (knapsack over their predicted marginals) for the
allocator-to-allocator comparison in Section~\ref{sec:allocation}. Aumann--Shapley scores are computed fit-free
by the path-integral estimator: the loss gradient is integrated by a 4-node midpoint rule along the
path that scales weights and unit-input activations linearly from full precision ($t=0$) to fully
quantized ($t=1$), one integral per demotion level ($16\to 8$ and $16\to 4$), with MoE routing frozen
along the path since the router's argmax is not differentiable. 
% This parameter-space path is the
% continuous relaxation of the multilinear-extension diagonal that Prop.~B.1 analyzes, where $t$ is
% instead the probability that each unit is quantized. The two integrals agree when the loss responds
% linearly to the per-unit perturbations and differ in general, and the correlation of 0.785 between the
% path-integral and fitted in-context coefficients on real models (App.~\ref{app:as}) is the empirical
% bridge between them.

\paragraph{Why joint least squares rather than per-coefficient sampling.} The plug-in estimate of
each Walsh coefficient from samples, $(1/N)\sum_n f(S_n)\chi_T(S_n)$ \citep{lmn1993}, is unbiased, but its
variance
carries the total second moment $\E_{\mup}[f^2]$, because every other coefficient's energy leaks in
through the
sampled, non-orthogonal design. Jointly regressing on all $L+1$ degree-$\le 1$ terms absorbs that
leakage into the fitted terms. Ordinary least squares is the classical control-variate regression
estimator, whose asymptotic variance is
only the energy of $f$ orthogonal to the fitted span, that is, the order-$\ge 2$ energy plus noise, rather
than $\Var(f)$ \citep{lavenberg1981cv,glynn2002cv,owen2013mcbook}. Random-design least squares is
stable and near-optimal already at $N \propto L\log L$ samples \citep{cohen2017optimalwls}. With
${\sim}90\%$ of the variance at order 1 (Section~\ref{sec:walsh}), regression gives roughly a $10\times$
per-coefficient
variance reduction at equal $N$, the reciprocal of the order-$\ge 2$ share, so the certificate also
quantifies the gain from regression. The same
regression-over-sampling argument recurs in Shapley estimation (KernelSHAP against permutation sampling,
\citealp{covert2021kernelshap}).

\paragraph{Statistical footing of the fits.} The coverage fit is ordinary nonlinear least squares:
consistent for the best coverage-family approximation of $f$ even when the model is misspecified
\citep{jennrich1969,white1981}, and the maximum-likelihood estimate under Gaussian noise when it is
correct. Two caveats: the objective is non-convex, so we initialize by linearizing through the
generator as described above, and on near-linear data the ceiling is weakly identified (only the
products $c\,a_i$ are pinned down), which perturbs $(c,a_i)$ individually but not predictions.

\paragraph{Measurement cost.} Cost is counted in measured configurations, since each is a forward pass
over the calibration set. The headline budget is $N\approx L$ configurations, against roughly $100L$
progressive-quantization evaluations for SPQE's permutation-sampled Shapley \citep{impq2025}.

\paragraph{Prediction benchmark in full.} Table~\ref{tab:master} reports the per-cell results of the
Section~\ref{sec:benchmark} benchmark at the small budget, and Fig.~\ref{fig:predn} the
error-versus-$N$ sweep on the 3B-wiki cell. The certificate column is the held-out residual of an
additive fit, the direct \thmone{} estimate: 300 configurations are sampled per cell from $\mu_{0.6}$,
an additive model is fit on the first 150, and the column reports the mean squared residual on the
remaining 150 divided by the damage variance. 

\begin{table}[t]
\centering\small
\caption{Prediction benchmark at the small-$N$ budget ($N=L/2$): Median held-out
relative error at half the $L$ single-layer measurements that isolate-and-sum requires. Entries are single-seed measurements over 120 held-out configurations.
Table~\ref{tab:seeds} reports 3-seed spreads on two W4A4 cells. Isolate-and-sum reaches 0.38--0.63 on the FP4-damaged models (0.6B, 4B, 30B),
against 0.08--0.19 for the better of coverage and additive on 0.6B and 4B. Two cells require comment.
On 8B the damage of the fully quantized model is small ($f(\text{all}) \approx 0.46$), all methods are
approximately tied, and isolate-and-sum looks adequate because there is
little to mis-price. The 30B MoE row is hard for every method (best 0.252). The last column is the
cell's order-$\ge 2$ share under $\mu_{0.6}$. It is
largest on the rows where every method's error is worst (4B-wiki at 0.21 and 30B at 0.30).}
\label{tab:master}
\resizebox{\textwidth}{!}{%
\begin{tabular}{llcccccccc}
\toprule
model & data & isolate-sum & additive & coverage & CoopQ & RF & GP & RBF & order-$\ge 2$ share \\
\midrule
Qwen3-0.6B      & wiki & 0.390 & \textbf{0.080} & 0.110 & 0.105 & 0.111 & 0.143 & 0.175 & .035 \\
Qwen3-0.6B      & code & 0.416 & 0.110 & 0.129 & \textbf{0.094} & 0.150 & 0.166 & 0.162 & .038 \\
Qwen2.5-3B      & wiki & 0.107 & 0.313 & \textbf{0.092} & 0.356 & 0.093 & 0.212 & 0.377 & .016 \\
Qwen2.5-3B      & code & 0.130 & 0.333 & 0.145 & 0.364 & \textbf{0.119} & 0.251 & 0.394 & .024 \\
Qwen3-4B        & wiki & 0.377 & 0.288 & \textbf{0.184} & 0.245 & 0.228 & 0.341 & 0.338 & .210 \\
Qwen3-4B        & code & 0.535 & 0.200 & 0.191 & \textbf{0.180} & 0.185 & 0.207 & 0.252 & .090 \\
Qwen3-8B        & wiki & \textbf{0.076} & 0.161 & 0.131 & 0.193 & 0.117 & 0.224 & 0.262 & .053 \\
Qwen3-8B        & code & 0.100 & 0.207 & 0.153 & 0.241 & \textbf{0.093} & 0.277 & 0.297 & .051 \\
Llama-3.2-3B    & wiki & 0.197 & 0.133 & \textbf{0.128} & 0.224 & 0.156 & 0.245 & 0.285 & .023 \\
Llama-3.2-3B    & code & 0.146 & 0.165 & \textbf{0.126} & 0.224 & 0.157 & 0.282 & 0.316 & .028 \\
Qwen3-30B (MoE) & wiki & 0.629 & 0.368 & 0.433 & 0.344 & 0.424 & 0.386 & \textbf{0.252} & .302 \\
\bottomrule
\end{tabular}}
\end{table}

\begin{figure}[t]
  \centering
  \includegraphics[width=0.72\linewidth]{fig3_predict_vs_N.png}
  \caption{Median held-out relative error vs.\ number of measured configurations on the 3B-wiki cell of
  Table~\ref{tab:master}. Coverage ($L+1$ parameters) sits at its error
  floor from $N=L/2$ (0.093, tied with the random forest at 0.094), the additive fit sits at 0.29--0.32
  for $N \le L$ and ties with coverage from $N \approx 4L$, and the Gaussian
  process needs roughly $4L$ configurations to match. From $N \approx L$ the random forest attains the
  lowest error on this cell, reaching 0.03 at $N \approx 8L$ against coverage's plateau at 0.08. The
  large-$N$ winner varies by cell (on 0.6B-wiki the Gaussian process leads and the forest trails every
  other method). Isolate-and-sum is constant at 0.11, since it has no fitted
  parameters, and its failures occur on the damaged cells of Table~\ref{tab:master}. CoopQ overfits at
  large $N$.}
  \label{fig:predn}
\end{figure}

\begin{table}[t]
\centering\small
\caption{Qwen-Image DiT (Section~\ref{sec:dit}): median held-out relative
error vs.\ $N$, $L=60$ blocks at per-block NVFP4. The final
column is the measured order-$\ge 2$
share.}
\label{tab:dit}
\begin{tabular}{cccccccc}
\toprule
$N$ & additive & coverage & pairwise & RBF-KR & RF & GP & order-$\ge 2$ share (cert.) \\
\midrule
30  & 0.356 & 0.151 & 0.425 & 0.378 & \textbf{0.142} & 0.368 & 0.36 \\
60  & 0.411 & 0.129 & 0.276 & 0.358 & \textbf{0.111} & 0.241 & \\
120 & 0.146 & \textbf{0.113} & 0.199 & 0.302 & 0.116 & 0.150 & \\
240 & 0.114 & 0.124 & 0.162 & 0.245 & \textbf{0.109} & 0.120 & \\
360 & \textbf{0.099} & 0.111 & 0.119 & 0.204 & 0.111 & 0.106 & \\
\bottomrule
\end{tabular}
\end{table}

\paragraph{Certificate estimators.} Table~\ref{tab:certest} reports the per-cell detail behind the
estimator comparison of Section~\ref{sec:tauval}: the root-mean-squared error, over 8 sampling seeds,
of estimating a lattice's true order-$\ge 2$ share from $N$ configurations sampled at $p=0.6$, for the
\propfive(i) forecast from the coverage fit and the held-out residual of an additive fit. The residual
estimator's agreement with the forecast, available at no extra measurement cost, is a per-instance
check of the coverage assumption itself.

\begin{table}[t]
  \centering
  \footnotesize
  \begin{tabular}{lccccccc}
    \toprule
    & & \multicolumn{2}{c}{$N=28$} & \multicolumn{2}{c}{$N=56$} & \multicolumn{2}{c}{$N=112$} \\
    \cmidrule(lr){3-4}\cmidrule(lr){5-6}\cmidrule(lr){7-8}
    cell & true share & fcst & resid & fcst & resid & fcst & resid \\
    \midrule
    Qwen3-0.6B FP4 per-block  & .008 & \textbf{.060} & .139 & \textbf{.067} & .078 & .073 & \textbf{.068} \\
    Qwen3-8B FP4 per-block    & .019 & \textbf{.010} & .366 & \textbf{.009} & .204 & \textbf{.006} & .168 \\
    Qwen3-0.6B INT8 per-tensor & .325 & \textbf{.083} & .362 & \textbf{.105} & .216 & \textbf{.110} & .240 \\
    Qwen3-0.6B FP4 per-tensor & .812 & \textbf{.195} & 1.23 & \textbf{.119} & .319 & \textbf{.067} & .538 \\
    Qwen3-8B FP4 per-tensor   & .922 & \textbf{.235} & 1.01 & \textbf{.164} & .492 & \textbf{.079} & .614 \\
    \bottomrule
  \end{tabular}
  \caption{Estimating a cell's order-$\ge 2$ share from $N$ sampled configurations. Each estimator receives only
  $N$ configurations drawn from $\mup$ at $p=0.6$, and the draw is repeated with 8 independent seeds.
  An entry is the root-mean-squared error of the 8 resulting estimates against the true share. The two
  estimators are the \propfive(i) forecast computed from a coverage fit (fcst) and the held-out residual
  of an additive fit (resid), both defined in Section~\ref{sec:tauval}.
  The forecast is best at 14 of 15 cell-budget pairs. 
  % The three per-tensor rows are lattices measured
  % for this comparison (10 rounding draws, 128-token calibration sequences). They are distinct from the
  % per-tensor certification cell of Table~\ref{tab:ordinal} (16 draws, 32-token calibration), whose
  % noise-corrected share is 0.67, so the share of a cell is quoted only with its own lattice.
  }
  \label{tab:certest}
\end{table}

\paragraph{Allocation.} Given a memory budget
$\mathrm{eff} = \sum \mathrm{numel}\cdot\mathrm{bits} / \sum \mathrm{numel}$ (kernel-independent, with
numel a tensor's element count), a
3-level $\{16,8,4\}$ knapsack minimizes predicted $\Delta\lce$ per effective bit. A unit is a weight tensor
or a fused group that must share one format, such as the q/k/v projections or one MoE layer's experts.
The served recipes
additionally tie fusion groups (q/k/v into one \code{qkv\_proj}, one expert format per MoE layer block) so
the per-layer-mixed checkpoint loads in TensorRT-LLM. Baselines are matched to the same $\mathrm{eff}$.
The two fit-free variants price demotions differently: AS-add uses the path-integral coefficients
directly, so predicted damage adds across units, while AS-cov converts the same coefficients to
break-rates, with the ceiling calibrated from one measured corner (the all-4-bit configuration), so
damage composes through the coverage form. The two structural rules use no damage measurements.
Naive-uniform lowers every unit together through the $\{16,8,4\}$ ladder and, if the resulting uniform
level falls below the target effective bits, restores the smallest tensors to 16 bits until the next
restoration would exceed the budget. Naive-routed places routed experts at the lowest level and protects
attention.

\input{generated/tab_ge4served30}
\input{generated/tab_ge4served235}

\paragraph{GPTQ ladder details (Section~\ref{sec:allocsub4}).} Grids are symmetric with one scale per
128 input columns, and each linear layer's Hessian is accumulated from the scoring half of the task
calibration set. The all-units GPTQ-4 calibration KL of 0.054 is comparable to NVFP4 RTN at 0.047, and
at 2 bits RTN gives 7.0 against GPTQ's 13.6, so 2 bits is unusable as a uniform level under either
quantizer. For the fused expert blocks, all experts
in a layer share the Hessian of the block's full token stream: the first projection's Hessian is
accumulated from the block input, and the second projection's from intermediate activations recomputed
on that input, a standard simplification that ignores the router's token-to-expert assignment. At the
235B scale the all-units GPTQ-3 calibration KL is 0.115.

\paragraph{Sub-4 served detail (Section~\ref{sec:allocsub4}).} Table~\ref{tab:sub4served30} gives the
full served grid on the reasoning benchmarks for every method and budget. The reasoning benchmarks
are AIME-2025 (30 problems), GPQA diamond (198 questions), and MMLU-Pro (a fixed
1{,}000-question subset), decoded at temperature 0.6, top-$p$ 0.95, top-$k$ 20, with a 32{,}768-token
budget on AIME-2025 and 16{,}384 elsewhere. The ModelOpt configuration scores at the chance level of
each multiple-choice format (23.2 on four-option GPQA, 9.0 on ten-option MMLU-Pro).

\input{generated/tab_sub4served30}



\paragraph{Joint measurement of the allocation tables.} 
% All cells of a given allocation table are
% produced by one pipeline on the same 16 evaluation sequences, so between-method comparisons are paired.
% The fake-quant evaluation is deterministic once the software environment is fixed (library versions are
% recorded alongside the result files), so every cell is exactly reproducible. All scoring signals use the
% scoring half of the calibration set, including the int8-outlier severity and the 300 fit configurations
% of \code{gp-alloc}, \code{rf-alloc}, and the fitted variants. 
CoopQ's quadratic program is solved by
greedy plus local search (SCIP unavailable). Its published ${\sim}100L$ budget (${\sim}24$k evaluations)
was not run (Section~\ref{sec:allocmoe} reports the 1{,}205- and 942-evaluation tiers). An ablation
isolates the cause of its deficit: the additive and coverage models fit on the
942 configurations CoopQ's permutation walks visited, allocating on its alphabet, land at
CoopQ's
level ($+.004$--$.008$, paired-significant), because configurations visited along a permutation walk are
efficient
for averaging marginal contributions but nearly collinear as regression data, while the same regression
on 300 independently sampled configurations lands within $.002$ of AS at every contended budget (within
$.001$ at eff $\ge 4.5$).

\paragraph{Hardware and software.} Scoring and fake-quant run on a single B200 (device-map
pipeline-parallel across $8\times$B200 for the 235B/355B/753B models), and serving runs on $1\times$B200
(30B, TP=1) up to $8\times$B200 (GLM, TP=8). PyTorch and Hugging Face Transformers handle fake-quant,
NVIDIA ModelOpt and TensorRT-LLM handle the $\ge 4$-bit
deployment recipes.
SGLang serves the sub-4 benchmarks from bf16-exported checkpoints and the GLM 5.2 benchmarks
(Section~\ref{sec:allocsub4}),
vLLM serves the Qwen3.6-27B comparison (Section~\ref{sec:alloc27b}), and
scikit-learn
provides the surrogate baselines.

\section{Interpretability of the fitted parameters}\label{sec:interp}

\paragraph{Budget pricing.} Fitting coverage on ${\sim}200$ configurations of Qwen3-0.6B (all 28
layers, activations quantized at FP4 per-block-32, configurations drawn at $p=0.6$), we price
unseen configurations of the form ``protect the $m$ most fragile layers and quantize the rest'',
comparing the predicted $\Delta L$, which costs no forward passes once the model is fit, against the
measured value (Table~\ref{tab:budget}). Across the deployment range, protecting a few to about half
the layers, coverage prices the configuration to within 4\%, whereas isolate-and-sum's prediction is
$2$--$3\times$ the truth. The lone breakdown is quantizing a single layer, where the model, fit on
configurations drawn at $p=0.6$, extrapolates far outside its fit distribution and misses by
${\sim}95\%$.

\begin{table}[t]
\centering\small
\caption{Budget pricing on unseen protect-$m$ configurations (Qwen3-0.6B).}
\label{tab:budget}
\begin{tabular}{cccc}
\toprule
protected $m$ & predicted $\Delta L$ & measured $\Delta L$ & rel.\ err. \\
\midrule
0 (all quantized) & 6.26 & 6.50 & 3.7\% \\
2  & 6.20 & 6.21 & 0.2\% \\
4  & 6.11 & 6.05 & 0.9\% \\
8  & 5.81 & 5.75 & 1.0\% \\
12 & 5.31 & 5.52 & 3.9\% \\
\bottomrule
\end{tabular}
\end{table}

\paragraph{The fragility scores are corpus-portable.} Fitting the break-rates $a_i$ independently on
wiki and on code yields strongly agreeing vectors (Spearman $\rho = 0.74$, Pearson $r = 0.86$ on
Qwen3-0.6B; $\rho = 0.78$, $r = 0.93$ on Qwen2.5-3B), with the same deep layers fragile in both
corpora. 
%The break-rates are therefore a property of the model rather than of the calibration text:
We propose to fit $a_i$ once on a source corpus, and to deploy on a new corpus refit only the single scalar ceiling
$c$. At small budgets this one-parameter transfer attains lower error than the $(L{+}1)$-parameter
in-domain fit, 0.058 against 0.098 at 9 configurations for code-to-wiki on both 0.6B and 3B, and 0.076
against 0.110 at 36 configurations for wiki-to-code on 3B.

\paragraph{Multi-seed robustness.} Table~\ref{tab:seeds} reports median held-out relative error over 3
seeds at the $N\approx L$ budget. Coverage attains the lowest error on 0.6B and ties the
Gaussian process on 3B, isolate-and-sum fails uniformly across seeds, and 3B's higher order-$\ge 2$
share (0.26 against 0.6B's 0.05) is exactly where coverage's margin over the additive fit is largest,
so the certificate forecasts where the saturation structure pays off before any predictor is fit.

\begin{table}[t]
\centering\small
\caption{Multi-seed robustness at the $N\approx L$ budget, with the certificate per row.}
\label{tab:seeds}
\resizebox{\textwidth}{!}{%
\begin{tabular}{lcccccccc}
\toprule
 & isolate-sum & additive & coverage & CoopQ & RF & GP & RBF & order-$\ge 2$ share (cert.) \\
\midrule
Qwen3-0.6B, $N{=}L{=}28$ & $0.474\pm0.010$ & $0.144\pm0.008$ & $\mathbf{0.085\pm0.013}$ & $0.113\pm0.002$ & $0.133\pm0.003$ & $0.099\pm0.002$ & $0.134\pm0.018$ & 0.05 \\
Qwen2.5-3B, $N{=}L{=}36$ & $2.798\pm0.030$ & $0.221\pm0.036$ & $\mathbf{0.098\pm0.009}$ & $0.143\pm0.019$ & $0.109\pm0.005$ & $0.095\pm0.008$ & $0.142\pm0.012$ & 0.26 \\
\bottomrule
\end{tabular}}
\end{table}

\section{Aumann--Shapley recovers the in-context coefficient at an effective density (proof under
coverage)}\label{app:as}

There are three per-layer slopes, all gradients of the saturating loss at different anchors: the
empty-anchor slope \citep[HAWQ, isolated;][]{hawq2019,dong2020hawqv2}, $s_i = f(\{i\}) = c\,a_i$, inflated
by the factor of Section~\ref{sec:anchor}; the deployment-point slope, which is the optimal additive coefficient,
$w_i = c\,a_i\prod_{j\ne i}(1 - p\,a_j)$, optimal but requiring a fit at the $p$-biased operating point;
and the path-average slope \citep[Aumann--Shapley;][]{aumann1974values,sundararajan2017ig}, which is
fit-free,
\[
  AS_i = c\,a_i\int_0^1 \prod_{j\ne i}(1 - t\,a_j)\,dt .
\]
Completeness, the property that the coefficients sum exactly to the total damage, is the classical
efficiency property of the Aumann--Shapley value and does not require the coverage form.
Under coverage it is immediate, since
$\frac{d}{dt}\bigl[c - c\prod_j(1 - t\,a_j)\bigr] = \sum_i c\,a_i\prod_{j\ne i}(1 - t\,a_j)$, and
integrating over $t \in [0,1]$ gives $\sum_i AS_i = f(\text{all})$ with no remainder (measured residual
on real models 0.06\%). To state what the path integral recovers, write
$w_i(q) = c\,a_i\prod_{j\ne i}(1 - q\,a_j)$ for the optimal additive coefficient when the deployment
density is $q$.
%  the corollary of \thmone{} and \thmfour{} in App.~\ref{app:proofs} applied with $q$ in
% place of $p$, so that the fitted coefficient above is $w_i = w_i(p)$.

\begin{customprop}{B.1}\label{prop:asrec}
Let $f$ be the coverage model with break-rates $a_1, \ldots, a_L \in [0,1)$ and let $p \in (0,1)$.
(i) For every layer $i$ there is an effective density $t_i^{*} \in (0, \tfrac12]$ such that
$AS_i = w_i(t_i^{*})$, so the fit-free path integral equals the optimal additive coefficient evaluated
at $t_i^{*}$ in place of $p$. (ii) For all layers $i$ and $k$, $AS_i > AS_k$ if and only if
$w_i > w_k$, if and only if $a_i > a_k$, so the path-average and deployment-point coefficients order
the layers identically at every deployment density.
\end{customprop}

\begin{proof}
(i) Write $g_i(t) = \prod_{j\ne i}(1 - t\,a_j)$, so that $AS_i = c\,a_i\int_0^1 g_i$ and
$w_i(q) = c\,a_i\,g_i(q)$. If $a_j = 0$ for every $j \ne i$ then $g_i \equiv 1$ and any $t_i^{*}$
serves. Otherwise $g_i$ is positive, continuous, strictly decreasing, and convex on $[0,1]$: with
$h_j(t) = -a_j/(1 - t\,a_j)$, differentiating the product gives $g_i' = g_i\sum_{j} h_j$ and
$g_i'' = g_i\sum_{j \ne l} h_j h_l \ge 0$, where the sums run over indices other than $i$ and the
$h_j$ agree in sign. Jensen's inequality for the uniform measure on $[0,1]$ gives
$\int_0^1 g_i \ge g_i(\tfrac12)$, strict monotonicity gives $\int_0^1 g_i < g_i(0)$, and the
intermediate value theorem then yields a unique $t_i^{*} \in (0, \tfrac12]$ with
$g_i(t_i^{*}) = \int_0^1 g_i$. Multiplying by $c\,a_i$ gives $AS_i = w_i(t_i^{*})$.
(ii) Factor out the full survival product $G(t) = \prod_{j}(1 - t\,a_j)$, positive on $[0,1]$ because
every $a_j < 1$. Then
\[
  w_i = c\,G(p)\,\frac{a_i}{1 - p\,a_i},
  \qquad
  AS_i = c\int_0^1 G(t)\,\frac{a_i}{1 - t\,a_i}\,dt ,
\]
and for each fixed $t \in [0,1]$ the map $x \mapsto x/(1 - t\,x)$ is strictly increasing on $[0,1)$.
Hence $a_i > a_k$ implies $w_i > w_k$ and $AS_i > AS_k$, while $a_i = a_k$ forces equality in both, and
the two orderings coincide with the ordering by break-rate.
\end{proof}

Part (i) places the fit-free
coefficient in the same family as the fitted one, evaluated at the effective density $t_i^{*}$ rather
than at the deployment density $p$ (since $t_i^{*} \le \tfrac12 < p$ at our $p=0.6$, the path-average
coefficient sits above the deployment-point one whenever other layers have nonzero break-rates). Part
(ii) shows that this difference in evaluation point never reorders layers. What remains numerical is
the magnitude agreement between $AS_i$ and
$w_i(p)$: the correlation is 0.9997 on coverage-generated data and 0.785 on real models, while the
empty-anchor $s_i$ is inflated by the factor ${\approx}\,e^{p\sum_j a_j}$. Since $s_i = c\,a_i$
is itself increasing in the break-rate, part (ii) extends to the empty anchor under the coverage model,
so the failure of isolated sensitivity documented in Section~\ref{sec:anchor} is a failure of magnitude rather
than of order. 

\section{Proofs (Theorems 3 and 4, Proposition 5, the two anchors, kernel eigenvalues)}\label{app:proofs}


\paragraph{Notation.} Under the product measure $\mup$ (each $x_i \sim \mathrm{Bernoulli}(p)$,
independent), the space of functions on $\{0,1\}^L$ carries the inner product
$\langle g,h\rangle = \E_{\mup}[gh]$. With the standardized coordinates
$\phi_i(x) := (x_i-p)/\sqrt{p(1-p)}$ ($\E[\phi_i]=0$, $\E[\phi_i^2]=1$), the Walsh functions
$\chi_T := \prod_{i\in T}\phi_i$ ($\chi_\emptyset := 1$) form an orthonormal basis: for $T\ne T'$,
independence factorizes $\E[\chi_T\chi_{T'}]$ per coordinate, and each coordinate in $T\triangle T'$
contributes $\E[\phi_i]=0$. Hence $f = \sum_T \fhat(T)\chi_T$ with $\fhat(T)=\langle f,\chi_T\rangle$,
$\E[f]=\fhat(\emptyset)$, and (Parseval) $\Var_p(f) = \sum_{T\ne\emptyset} \fhat(T)^2$
\citep{odonnell2014boolean}.

\paragraph{Proof of Theorem 3 (M\"obius spectrum).} For $T\ne\emptyset$,
$\varphi(T) = \sum_{R\subseteq T}(-1)^{|T|-|R|} f(R)$ with $f(R) = c - c\prod_{i\in R}(1-a_i)$. The
constant $c$ alternates to zero over the subsets of a nonempty $T$. For the product part, the alternating
subset-sum of a product factorizes coordinatewise,
\[
\sum_{R\subseteq T}(-1)^{|T|-|R|}\prod_{i\in R}g_i = \prod_{i\in T}(g_i-1).
\]
With $g_i = 1-a_i$ this gives
$-c\prod_{i\in T}(-a_i) = c\,(-1)^{|T|+1}\prod_{i\in T}a_i$. \qed

\paragraph{Proof of Theorem 4 (Walsh spectrum).} Write $f = c - H$ with the headroom
$H(x) = c\prod_i g_i(x_i)$, $g_i(x_i) = 1 - a_i x_i$. Constants are orthogonal to every $\chi_T$ with
$T\ne\emptyset$, so $\fhat(T) = -\hat H(T)$ for $T\ne\emptyset$, while
$\fhat(\emptyset) = \E[f] = c - \hat H(\emptyset) = c\bigl(1-\prod_i(1-a_ip)\bigr)$. Because $H$ and $\chi_T$ are both products over
\emph{independent} coordinates, the expectation factorizes:
$\hat H(T) = c\prod_{i\in T}\E[g_i\phi_i]\cdot\prod_{i\notin T}\E[g_i]$. The one-coordinate integrals
(using $x_i^2=x_i$): $\E[g_i] = 1-a_ip$ and
$\E[g_i\phi_i] = -a_i\,\E[x_i(x_i-p)]/\sqrt{p(1-p)} = -a_i\sqrt{p(1-p)}$. Substituting and negating gives
$\fhat(T) = c\,(-1)^{|T|+1}(p(1-p))^{|T|/2}\prod_{i\in T}a_i\cdot\prod_{i\notin T}(1-a_ip)$. \qed



\begin{proof}[Proof of \propfive(ii) and (iii)]
Notation as in Section~\ref{sec:tau}, with $B := \tau^2 = \sum_i\beta_i$. By \propfive(i),
$\sharegeq = S_2/S_1$ with $S_1 := \sum_{k\ge 1}\es{k} = \prod_i(1+\beta_i)-1$ and
$S_2 := \sum_{k\ge 2}\es{k}$.

\emph{(ii)} The leading term is exact algebra: $(\sum\beta_i)^2 = \sum\beta_i^2 + 2\es{2}$, so
$\es{2} = \tfrac12(B^2 - \sum\beta_i^2) = \tfrac12 B^2(1 - 1/\Leff)$, hence
$\es{2}/\es{1} = \tfrac12 B(1 - 1/\Leff)$. For the error we use $k!\,\es{k} \le B^k$ (expand
$B^k = (\sum\beta)^k$: the $k!\,\es{k}$ injective monomials are a subset of its nonnegative terms), whence
$S_2 - \es{2} = \sum_{k\ge 3}\es{k} \le \sum_{k\ge 3}B^k/k! \le (B^3/6)e^B$ (termwise,
$3!(k-3)! \le k!$) and $S_1 \le e^B - 1$. \emph{Upper:}
$\sharegeq \le S_2/\es{1} \le \es{2}/B + (B^2/6)e^B$. \emph{Lower:}
$\sharegeq \ge \es{2}/S_1 \ge \es{2}/(e^B-1) = (\es{2}/B)\cdot B/(e^B-1) \ge (\es{2}/B)(1 - B/2)
\ge \es{2}/B - B^2/4$, using the elementary inequality $x/(e^x-1) \ge 1 - x/2$ for $x \ge 0$ (the power
series of $x - (e^x-1)(1-x/2)$ is $\sum_{k\ge 3}\frac{k-2}{2\,k!}x^k$, with nonnegative coefficients) and
$\es{2}/B \le B/2$. Both one-sided deviations are $\le (B^2/4)e^B$, which is the stated remainder bound.
The profile's third moment (via $\es{3}$) first appears at $O(\tau^4)$.

\emph{(iii)} At fixed $B$, $\sharegeq = 1 - B/(P-1)$ is strictly increasing in
$P := \prod_i(1+\beta_i)$, and $\log P = \sum_i \log(1+\beta_i)$ is a symmetric sum of one concave scalar
function over the coordinates, the canonical Schur-concave form \citep{marshall2011majorization}, so
$\sharegeq$ is Schur-concave at fixed $B$. The uniform profile $\beta_i \equiv B/L$ is majorized by
every profile of total $B$, hence is the maximizer (equivalently, Jensen on the strictly concave
$\log(1+\cdot)$ gives $\log P \le L\log(1+B/L)$, with equality iff all $\beta_i$ are equal). The outer
bound follows since $L \mapsto (1+B/L)^L$ increases to $e^B$ (AM--GM on $L$ copies of $1+B/L$ and one
copy of $1$). As $\Leff \to 1$ the share vanishes, since a one-fragile-layer $f \approx c\,a_i\,x_i$
is a one-coordinate, hence exactly degree-1, function.
\end{proof}

\paragraph{Corollary (the two anchors of Section~\ref{sec:anchor}).} By Theorem 1 the optimal degree-$\le 1$
predictor is the projection $\fhat(\emptyset) + \sum_i \fhat(\{i\})\chi_{\{i\}}$. Converting
$\chi_{\{i\}}$ back to the $\{0,1\}$ coordinate gives the optimal additive coefficient
$w_i = \fhat(\{i\})/\sqrt{p(1-p)} = \E_p[f\mid x_i=1] - \E_p[f\mid x_i=0]$, the in-context marginal, which
under coverage (Theorem~4 with $T=\{i\}$) is $w_i = c\,a_i\prod_{j\ne i}(1-a_jp)$. The empty-anchor
slope is $s_i = f(\{i\}) = c\,a_i$ directly from the coverage model. Hence the inflation factor
$s_i/w_i = 1/\prod_{j\ne i}(1-a_jp) \approx e^{p\sum_j a_j}$ quoted in Section~\ref{sec:anchor}. \qed

% The statements are numerically verified. The product form matches brute-force Theorem-4 enumeration to
% ${\sim}10^{-16}$ over random $(L, a, c, p)$, the $(\tau^4/4)e^{\tau^2}$ bound and the symmetric upper
% bound of \propfive(iii) hold on
% 5{,}000 random heterogeneous profiles ($L=2\ldots 40$, zero violations, worst $|R|$/bound $= 0.31$), and
% $|R|/\tau^4$ is stable as a fixed profile is scaled down, confirming the error order.

\paragraph{The symmetric case.} For completeness, the equality case of \propfive(iii) with its order
spectrum.

\begin{corollary*}[symmetric saturation law]
If $a_i \equiv a$, then $\beta_i \equiv \beta = p(1-p)\,a^2/(1-ap)^2$ and $\tau^2 = L\beta$, the order
spectrum is binomial,
\[
  W_k = c^2\binom{L}{k}\,(p(1-p))^k a^{2k}(1-ap)^{2(L-k)},
\]
and the certificate collapses to the one-parameter law
\[
  \sharegeq = 1 - \frac{\tau^2}{(1+\tau^2/L)^L - 1}
  \;\xrightarrow{\;\tau\to 0\;}\; \frac{L-1}{2L}\,\tau^2 \;\approx\; \tfrac{1}{2}\tau^2 .
\]
\end{corollary*}

\begin{proof}
With equal break-rates every size-$k$ subset contributes the same product, so
$\es{k} = \binom{L}{k}\beta^k$ and $\prod_i(1+\beta_i) = (1+\beta)^L$, and \propfive(i) reduces to the
stated law, which is also the equality case of \propfive(iii). The $W_k$ formula is \thmfour{} with
equal $a_i$ times the $\binom{L}{k}$ count of size-$k$ sets, and the small-$\tau$ limit is
\propfive(ii) with $\Leff = L$.
\end{proof}

\paragraph{Walsh eigenvalues of the RBF kernel (Section~\ref{sec:ood}).} On binary vectors the squared
Euclidean distance collapses to the size of the symmetric difference, $\lVert x - x'\rVert^2 =
\sum_i (x_i - x_i')^2 = |x \oplus x'|$, so the RBF kernel is a function of the XOR alone,
$K(x,x') = \kappa(x \oplus x')$ with $\kappa(u) = e^{-\gamma \sum_i u_i}$. For any such kernel, the
operator $(T_K f)(x) = \E_{x'}[K(x,x')f(x')]$ under the uniform measure ($p=\tfrac12$) is diagonal in
the Walsh basis. Fix $x$ and reparametrize the average over $x'$ by $u := x \oplus x'$. Hence
\[
  (T_K\chi_T)(x) = \E_{x'}\bigl[\kappa(x \oplus x')\,\chi_T(x')\bigr]
  = \E_u\bigl[\kappa(u)\,\chi_T(x \oplus u)\bigr].
\]
The character identity
$\chi_T(x \oplus u) = \chi_T(x)\,(-1)^{\sum_{i\in T} u_i}$ then gives
$T_K\chi_T = \lambda_T\,\chi_T$ with $\lambda_T = \E_u\bigl[\kappa(u)\,(-1)^{\sum_{i\in T} u_i}\bigr]$.
This expectation factorizes over the independent uniform bits,
\[
  \lambda_T \;=\; \prod_{i\in T}\E\bigl[e^{-\gamma u_i}(-1)^{u_i}\bigr]
  \prod_{i\notin T}\E\bigl[e^{-\gamma u_i}\bigr]
  \;=\; \Bigl(\tfrac{1-e^{-\gamma}}{2}\Bigr)^{|T|}\Bigl(\tfrac{1+e^{-\gamma}}{2}\Bigr)^{L-|T|}
  \;=\; \Bigl(\tfrac{1+e^{-\gamma}}{2}\Bigr)^{L}\,\tanh(\gamma/2)^{|T|}.
\]
The eigenvalues are positive, depend on $T$ only through $|T|$, and decay geometrically in $|T|$ with
ratio $\tanh(\gamma/2) < 1$, which is the high-order shrinkage quoted in Section~\ref{sec:ood}. At
$p \ne \tfrac12$ the kernel is still a function of $x \oplus x'$, but the $\mup$-Walsh basis is no
longer exactly its eigenbasis, which is why Section~\ref{sec:ood} states the diagonalization as
approximate at $p = 0.6$.

\section{The log-headroom knapsack, and the generator the data selects}\label{app:knapsack}

\paragraph{Knapsack certificate.} Because coverage is linear in $b$, a service-level constraint of the
form
$\Delta\lce \le \varepsilon$ maps to a scalar budget $\sum_{i\in S} b_i \le B$ with
$B = -\log(1 - \varepsilon/c)$. The cheapest configuration meeting the constraint protects the set $P$
minimizing protection
cost subject to $\sum_{i\notin P} b_i \le B$ is a 0/1 knapsack, solvable exactly by
dynamic programming. The deliverable is a predicted $\Delta\lce$ for the chosen configuration,
quotable before committing ($\le 4\%$ error, App.~\ref{sec:interp}), which a ranking structurally cannot
produce and which a Gaussian process prices with $1.6$--$2.9\times$ coverage's error near the heavily
quantized corner where optimal configurations live (Section~\ref{sec:ood}).

\paragraph{Which generator the data selects.} The ceiling $c$ and membership in the class
$g\bigl(\sum_i b_i\bigr)$ are established empirically, by the marginal-damage measurements of
Section~\ref{sec:coverage} and the ordinal certification of Section~\ref{sec:ordinal}, while the
exponential generator (coverage/noisy-OR), which is what makes the headroom consumption multiplicative, is
justified only by
fit quality. Refitting the lattices with candidate strictly increasing generators (all with $L+1$
parameters, $p=0.6$) shows that prediction is insensitive to the choice (held-out relative error
${\approx}\,0.064$--$0.071$ for all), but the slope calibration, the ratio
$w_{\text{model}}/w_{\text{data}}$ of model-implied to measured in-context slopes, selects one: across
the five non-degenerate of six (model, corpus) cells, the probit-$\sqrt{\cdot}$ ratio is 0.93--1.03,
while the exponential ratio spans 0.86 down to ${\approx}0$ (exp, probit, tanh, and rational sit at
0.58--0.67 on Qwen3-0.6B). The saturation rule
best supported by the data is therefore the probit-$\sqrt{\cdot}$ generator
$g(u) = c\,\operatorname{erfc}\bigl(1/\sqrt{2u}\bigr)$, the probability that zero-mean Gaussian
corruption of variance $u$ exceeds a fixed threshold, scaled by the ceiling; the exponential coverage
model is its tractable approximation. 

% Two further measurements agree. The coverage-form
% Aumann--Shapley coefficients land at about half the fitted in-context coefficients under the
% exponential generator, with the ratio restored under probit-$\sqrt{\cdot}$. And the coverage-predicted
% anchor-inflation ratios of Section~\ref{sec:anchor} sit uniformly below the measured ones while
% correlating with them per layer within each cell (Pearson $0.40$--$0.92$ across the four per-block
% certification cells), the direction expected if the exponential generator compresses the slopes that
% the probit-$\sqrt{\cdot}$ generator matches.

\section{Extended related work}\label{app:related}

\paragraph{MoE-specific mixed precision.} For mixture-of-experts models, MC-MoE
\citep{mcmoe2025} solves a linear program over expert routing and frequency importance. It allocates by
an additive importance score and does not analyze the order of the damage or report a validity
certificate. The head-to-head comparisons against CoopQ and MC-MoE appear in
Section~\ref{sec:allocation}.

\paragraph{Path-integral sensitivity.} Integrated Gradients \citep{sundararajan2017ig} and the
Aumann--Shapley value \citep{aumann1974values} are the classical methods. PQI/ReQuant \citep{hu2025pqi}
applies the path integral to weight-level sensitivity but assumes per-weight effects sum, and QIG
\citep{qig2026} is
token-level for vision-language models. We apply the path integral to the layer-selection set function
(App.~\ref{app:as}).

\paragraph{Rotation methods.} QuaRot, SpinQuant, and QuIP\#
\citep{ashkboos2024quarot,liu2024spinquant,tseng2024quipsharp} flatten the representation before
quantizing rather than allocating precision. They alter the base quantizer on which any allocator
operates, so they compose with allocation (Section~\ref{sec:future}). A rotation that flattens per-layer
heterogeneity also removes the per-layer differences that allocation exploits.

\paragraph{Low-order surrogates of subset functions outside quantization.} The analytical approach has
precedents in several adjacent literatures: functional ANOVA over configuration spaces
\citep{hutter2014fanova}; Datamodels \citep{ilyas2022datamodels}, whose finding that linear surrogates of
training-data-subset effects work well is the data-subset analog of our first-order finding; SPEX
\citep{spex2025}, sparse Boolean-Fourier surrogates of LLM outputs over input-masking subsets; the
Harsanyi-interaction line \citep{ren2023harsanyi}, which applies M\"obius machinery to input coalitions;
and Walsh surrogates for pseudo-Boolean black-box optimization \citep{lepretre2019walsh,bocs2018}. 

\paragraph{In-context and isolated measurements.} That an effect measured in isolation differs from the effect in context
is known qualitatively from coalition-value pruning \citep{ancona2020shapley}, the joint-versus-sum
analysis of Combinatorial Brain Surgeon for weight pruning \citep{cbs2022}, the Hydra-effect self-repair
mechanism \citep{mcgrath2023hydra}, the permutation-sampled marginals of CoopQ's SPQE estimator
\citep{impq2025}, and the
layer-subset surrogates now used for LLM pruning \citep{pruninggame2026}, all of which work around the
same anchor problem. What we contribute is the application to mixed-precision quantization,
specifically, the closed-form inflation factor under saturation, its measured
magnitude on modern LLMs, and a fit-free estimator of the correct coefficient (Section~\ref{sec:predictor}), one
Aumann--Shapley path integral in place of the SPQE estimator's ${\sim}100L$ quantized evaluations.
\begin{table}[t]
\centering\footnotesize
\caption{Sub-4 allocation on the GPTQ ladder (Qwen3-30B), calibration KL on the evaluation-task
corpus (96 MBPP, 64 gsm8k-train, 64 MATH, 64 MMLU-Pro validation prompts; evaluation half, 16
shared sequences). Every method is derived and measured on this corpus. $*$ marks methods
significantly different from the additive allocator under a paired $t$-test over the 16 shared
sequences, Holm-corrected within each budget row at the 5\% level. Parentheticals give each
method's measurement cost: \emph{meas.}\ counts isolated per-unit damage measurements on the
ladder (additive prices each unit's $4\to3$ and $3\to2$ steps; isolated-rank ranks by the same
damages), \emph{cfg.}\ counts sampled mixed configurations the fitted, GP, and RF predictors are
fit on, and CoopQ's count is its permutation-walk evaluations.}
\label{tab:sub4kl}
\resizebox{\textwidth}{!}{%
\begin{tabular}{ccccccccccc}
\toprule
eff & additive & CoopQ & fitted & gp & rf & isolated-rank & uniform & ModelOpt & MC-MoE & int8-outlier \\
 & (96 meas.) & (376 meas.) & (40 cfg.) & (40 cfg.) & (40 cfg.) & (192 meas.) & & & & \\
\midrule
3.5 & 0.086 & 0.092 & 0.161$^*$ & 0.178$^*$ & 0.176$^*$ & \textbf{0.082}$^*$ & 0.099$^*$ & 0.300$^*$ & 0.271$^*$ & 0.386$^*$ \\
3.25 & \textbf{0.107} & 0.110 & 0.214$^*$ & 0.223$^*$ & 0.287$^*$ & 0.107 & 0.124$^*$ & 0.620$^*$ & 0.528$^*$ & 0.703$^*$ \\
3.0 & \textbf{0.145} & \textbf{0.145} & 0.293$^*$ & 0.337$^*$ & 0.449$^*$ & \textbf{0.145} & \textbf{0.145} & 0.952$^*$ & 1.098$^*$ & 1.110$^*$ \\
2.9 & 0.215 & 0.228$^*$ & 0.322$^*$ & 0.350$^*$ & 0.548$^*$ & \textbf{0.191}$^*$ & 0.236$^*$ & 1.042$^*$ & 1.222$^*$ & 1.229$^*$ \\
2.75 & 0.415 & 0.354$^*$ & 0.422 & 0.412 & 0.849$^*$ & 0.389 & \textbf{0.319}$^*$ & 1.255$^*$ & 1.366$^*$ & 1.370$^*$ \\
2.5 & 0.854 & \textbf{0.631}$^*$ & 0.715$^*$ & 0.710$^*$ & 1.445$^*$ & 1.026$^*$ & 0.798$^*$ & 1.605$^*$ & 1.695$^*$ & 1.736$^*$ \\
2.3 & 1.563 & \textbf{1.122}$^*$ & 1.298$^*$ & 1.349$^*$ & 2.121$^*$ & 1.627$^*$ & 1.649$^*$ & 2.079$^*$ & 2.215$^*$ & 2.215$^*$ \\
2.1 & 2.482 & \textbf{2.285}$^*$ & 2.472 & 2.472 & 2.984$^*$ & 2.482 & 2.617$^*$ & 2.967$^*$ & 3.057$^*$ & 2.967$^*$ \\
\bottomrule
\end{tabular}}
\end{table}

\begin{table}[t]
\centering\footnotesize
\caption{Qwen3-235B below four bits on the evaluation-task corpus. Top: calibration KL of the
three directly constructed allocations (same ladder and corpus protocol as
Table~\ref{tab:sub4kl}). Bottom: served accuracy, AIME-2025 (8 seeds) / GPQA diamond /
MMLU-Pro, for the five listed configurations (GPTQ weight-only checkpoints); bf16 anchors 80.0/71.7/82.9.
The additive allocation preserves reasoning at 3.5 effective bits and degrades through 3.0;
every earlier-derived method collapses by 2.9. The rows marked eval.-task deriv.\ are
allocations derived on the evaluation-task corpus itself. 
Both remain functional at 2.9 (42.9 and 41.2 AIME with no
generation reaching the token cap) The surviving configurations are those whose calibration KL (0.44) sits below the
termination threshold of Section~\ref{sec:allocsub4}.}
\label{tab:sub4b235}
\begin{tabular}{lcccc}
\toprule
calibration KL & eff 3.5 & eff 3.0 & eff 2.9 & eff 2.75 \\
\midrule
additive & 0.167 & 0.292 & 0.435 & 1.162 \\
uniform & 0.217 & 0.297 & 2.006 & 4.674 \\
ModelOpt & 1.477 & 5.642 & 6.274 & 7.462 \\
\midrule
served & & & & \\
additive & 73.3/68.7/81.8 & 39.6/52.0/69.6 & 2.1/31.8/41.0 & 0.0/22.2/10.0 \\
additive (eval.-task deriv.) & 70.4/66.7/81.4 & 58.3/60.1/77.1 & 42.9/48.0/66.8 & 0.0/18.7/11.6 \\
coverage (eval.-task deriv.) & 71.7/66.7/80.9 & 59.6/62.6/76.3 & 41.2/44.4/66.7 & 0.0/25.3/10.9 \\
uniform & 66.2/65.7/80.6 & 60.4/61.1/77.2 & 5.0/16.7/21.3 & 0.0/0.0/0.2 \\
ModelOpt & 0.0/18.2/8.7 & 0.0/2.0/1.0 & 0.0/1.0/0.3 & 0.0/0.0/0.2 \\
\bottomrule
\end{tabular}
\end{table}

\begin{table}[t]
\centering\footnotesize
\caption{Served accuracy of the Qwen3-235B allocations at effective bits 4.2--8.0 over the
$\{16,8,4\}$ levels (allocations derived on the evaluation-task corpus of
Table~\ref{tab:sub4kl}).
Each cell gives AIME-2025 (average of 8 sampled seeds) / GPQA diamond / MMLU-Pro; bf16 baselines
80.0/71.7/82.9. Every measured cell is within seed noise of the anchor on all three benchmarks.}
\label{tab:ge4served235}
\begin{tabular}{lccc}
\toprule
 & eff 4.2 & eff 5.0 & eff 8.0 \\
\midrule
uniform & 79.6/70.2/82.6 & 78.3/69.2/82.3 & 81.2/70.7/82.7 \\
AS-additive & 77.1/70.7/83.4 & 80.4/69.2/83.0 & 82.1/68.2/83.2 \\
ModelOpt & 78.3/69.2/82.5 & 80.4/73.2/81.8 & 79.6/67.2/83.7 \\
\bottomrule
\end{tabular}
\end{table}

\begin{table}[t]
\centering\footnotesize
\caption{Served accuracy of every sub-4 Qwen3-30B configuration on the reasoning benchmarks.
Each cell gives AIME-2025 (average of 8 sampled seeds) / GPQA diamond / MMLU-Pro pass rates for
the method's configuration at that effective-bit budget All allocations are derived from damages measured on the
evaluation-task calibration corpus of Table~\ref{tab:sub4kl}. bf16 baselines: 73.3/61.6/76.2.}
\label{tab:sub4served30}
\begin{tabular}{lcccc}
\toprule
 & eff 3.5 & eff 3.0 & eff 2.9 & eff 2.75 \\
\midrule
additive & 68.8/62.1/73.7 & 59.6/51.0/73.4 & 49.6/44.4/67.1 & 23.3/31.3/54.3 \\
coverage & 65.4/56.1/74.6 & 61.7/53.0/74.0 & 55.0/47.0/68.1 & 25.4/34.3/55.4 \\
CoopQ & 66.2/55.6/74.1 & 59.2/53.5/72.1 & 46.2/47.5/67.3 & 14.2/27.3/55.2 \\
isolated-rank & 63.3/55.6/74.7 & 57.1/50.5/73.1 & 56.2/52.0/67.0 & 26.2/30.8/54.4 \\
uniform & 65.8/57.1/74.5 & 59.2/54.0/73.5 & 41.2/47.5/65.2 & 30.0/42.9/58.6 \\
fitted & 56.7/54.0/71.4 & 34.2/38.9/60.5 & 33.3/36.4/57.5 & 15.4/25.8/48.3 \\
gp & 53.3/49.5/69.3 & 31.2/32.8/58.1 & 30.0/37.9/55.1 & 3.3/25.8/45.4 \\
rf & 57.1/50.5/70.6 & 3.3/20.7/45.0 & 0.0/9.6/20.9 & 0.0/7.6/4.3 \\
ModelOpt & 40.0/31.3/60.4 & 0.0/18.7/8.7 & 0.0/22.7/9.1 & 0.0/18.2/6.3 \\
MC-MoE & 38.8/45.5/66.1 & 0.0/18.7/9.0 & 0.0/16.7/8.0 & 0.0/16.7/4.7 \\
int8-outlier & 35.8/31.8/53.8 & 0.0/18.7/7.0 & 0.0/13.1/--- & 0.0/12.1/4.0 \\
\bottomrule
\end{tabular}
\end{table}

\begin{table}[t]
\centering\small
\caption{Served AIME-2025 accuracy (average of 8 sampled seeds) of every Qwen3-30B allocation at
effective bits 4.2--8.0 over the $\{16,8,4\}$ levels,
served by SGLang; bf16 anchor 73.3. GPQA diamond spans 59.1--63.1 and MMLU-Pro 76.2--78.3 across the same
grid, no method ordering is significant at any budget (a
16-seed replication of the AS-additive--uniform gap at eff 4.5 gives 73.3
against 69.8, about 1.4 standard errors).}
\label{tab:ge4served30}
\begin{tabular}{lccccc}
\toprule
 & eff 4.2 & eff 4.5 & eff 5.0 & eff 6.0 & eff 8.0 \\
\midrule
AS-additive & 68.8 & 75.4 & 69.2 & 71.7 & 70.4 \\
AS-coverage & 69.2 & 69.6 & 69.6 & 71.2 & 70.4 \\
uniform & 68.8 & 69.6 & 69.2 & 69.6 & 72.1 \\
routed & 71.2 & 70.8 & 69.6 & 71.7 & 70.8 \\
ModelOpt & 70.0 & 70.0 & 72.9 & 70.8 & 68.8 \\
fitted & 70.8 & 71.7 & 71.2 & 73.8 & 69.6 \\
gp & 70.8 & 69.6 & 73.8 & 69.6 & 72.1 \\
rf & 72.5 & 73.3 & 70.8 & 69.2 & 72.5 \\
CoopQ & 67.5 & 75.8 & 69.6 & 68.8 & 74.2 \\
MC-MoE & 67.1 & 67.5 & 70.4 & 68.8 & 71.7 \\
int8-outlier & 71.7 & 67.1 & 72.5 & 73.3 & 75.0 \\
\bottomrule
\end{tabular}
\end{table}

% ---------------------------------------------------------------------------
% Theorem machinery. The numbering follows order of appearance (Thm 1, Thm 2 =
% the ordinal theorem, Thm 3 = Mobius spectrum, Thm 4 = Walsh spectrum,
% Prop 5 = heterogeneous certificate) and is load-bearing in cross-references,
% so theorem numbers are set manually via the custom* environments below.
% ---------------------------------------------------------------------------
\theoremstyle{plain}
\newtheorem{innercustomthm}{Theorem}
\newenvironment{customthm}[1]
  {\renewcommand\theinnercustomthm{#1}\innercustomthm}
  {\endinnercustomthm}
\newtheorem{innercustomprop}{Proposition}
\newenvironment{customprop}[1]
  {\renewcommand\theinnercustomprop{#1}\innercustomprop}
  {\endinnercustomprop}
\newtheorem{innercustomcor}{Corollary}
\newenvironment{customcor}[1]
  {\renewcommand\theinnercustomcor{#1}\innercustomcor}
  {\endinnercustomcor}
\newtheorem*{corollary*}{Corollary}
\newtheorem*{proposition*}{Proposition}

% Manual theorem-number shorthands used in prose.
\newcommand{\thmone}{Theorem~1}
\newcommand{\thmtwo}{Theorem~2}    % the ordinal theorem
\newcommand{\thmthree}{Theorem~3}  % Mobius spectrum
\newcommand{\thmfour}{Theorem~4}   % Walsh spectrum
\newcommand{\propfive}{Proposition~5}  % heterogeneous certificate
% Legacy aliases (kept so stray references still compile; prefer the new names).
\newcommand{\thmthreep}{\thmfour}
\newcommand{\propthreepp}{\propfive}

% ---------------------------------------------------------------------------
% Math shorthands.
% ---------------------------------------------------------------------------
\newcommand{\E}{\mathbb{E}}
\DeclareMathOperator{\Var}{Var}
\newcommand{\lce}{L_{\mathrm{CE}}}
\newcommand{\fhat}{\hat f}
\newcommand{\mup}{\mu_p}
\newcommand{\sharegeq}{\mathrm{share}_{\ge 2}}
\newcommand{\Leff}{L_{\mathrm{eff}}}
\newcommand{\es}[1]{e_{#1}}            % elementary symmetric polynomial

% Inline code / file names.
\newcommand{\code}[1]{\texttt{#1}}

\documentclass[11pt]{article}

% ---------------------------------------------------------------------------
% Preprint layout (arXiv-style, no conference class). Swap the class + drop
% geometry to move to a venue template later; the body lives in sections/.
% ---------------------------------------------------------------------------
\usepackage[margin=1in]{geometry}
\usepackage{amsmath,amssymb,amsthm,mathtools}
\usepackage{booktabs}
\usepackage{graphicx}
\usepackage{subcaption}
\usepackage{microtype}
\usepackage{natbib}            % authoryear; \citep / \citet
\usepackage{enumitem}
\setlist{topsep=2pt,itemsep=2pt,leftmargin=1.4em}
\usepackage{xcolor}
\usepackage{authblk}           % author/affiliation blocks with superscript labels
\usepackage[colorlinks=true,linkcolor=blue!55!black,citecolor=blue!55!black,urlcolor=blue!55!black]{hyperref}

\graphicspath{{figures/}{generated/}}

\input{macros}

\title{Saturation Makes Quantization Error Additive: A Coverage Model with a Certificate}

\author{Joshua Hill\thanks{Work done during an internship at Baseten Labs, Inc.}\\
Baseten Labs, Inc.\\  
David R.\ Cheriton School of Computer Science, University of Waterloo\\
}
\date{}

\begin{document}
\maketitle

\begin{abstract}
Mixed-precision quantization must decide which parts of a model to keep at higher precision. A common
premise, shared by sensitivity-based methods such as HAWQ and CoopQ, is that the loss from quantizing a set
of layers can be reconstructed from per-layer or pairwise sensitivities measured in isolation. We test this
premise at the 4-bit weight-and-activation precisions now being deployed, treating the change in loss
$f(S)$ from quantizing a layer set $S$ as a set function on the Boolean cube and
analyzing it through two classical changes of basis. This analysis yields two findings. First, across
configurations drawn from the deployment distribution, 85--93\% of the variance of $f$ is explained by
per-layer effects alone. Second, a monotone transform of a sum of per-layer terms reproduces $f$'s ranking of
configurations, misordering at most 2\% of pairs.

Both findings are consequences of the \emph{saturation} of the loss. Because the loss is bounded, each quantized
layer consumes a fraction $a_i$ (its break-rate) of the remaining headroom toward a ceiling $c$; the
resulting coverage model $f(S)=c\bigl(1-\prod_{i\in S}(1-a_i)\bigr)$ reproduces the measured variance
profile of $f$ to within a few percent from its $L$ fitted break-rates. This structure supports two
predictors of a configuration's loss, each with $L+1$ parameters. The additive model is the optimal
first-order predictor. By Parseval's identity its mean-squared error
equals the variance of $f$ left unexplained by per-layer effects, which we measure on full lattices,
estimate out of sample at full-network scale, and report with every result as
a certificate
of how well any additive model can do. The coverage model itself is the second predictor, and the two
predict comparably.

Saturation makes a layer's marginal damage depend on how many other layers are quantized, so sensitivities
measured in isolation, including a Hessian-trace HAWQ-v2, mis-price a configuration's loss by
71--376\%. An Aumann--Shapley path integral recovers the optimal additive coefficient up to a
layer-specific effective deployment density, preserving its ranking of layers exactly, a recovery we
prove under the coverage model. Both predictors price a configuration before deployment, match black-box
surrogates at small measurement budgets, degrade at most $2\times$ beyond the
configurations fit on, where the surrogates degrade $2$--$5\times$. As allocators at matched memory, they
attain the lowest KL divergence among the compared allocators at the contended budgets of
mixture-of-experts models from 30B to 355B parameters (17--21\% lower KL
divergence than the production ModelOpt baseline on Qwen3-30B).
Below four bits, the resulting allocations continue to solve code and reasoning tasks
at budgets where allocations from gradient sensitivities no longer produce terminating generations.
\end{abstract}

\input{sections/01_intro}
\input{sections/02_related}
\input{sections/04_walsh}
\input{sections/05_theory}
\input{sections/06_anchor}
\input{sections/07_predictor}
\input{sections/08_allocation}
\input{sections/09_scope}
\input{sections/11_future}
\input{sections/12_conclusion}

\bibliographystyle{plainnat}
\bibliography{references}

\clearpage
\appendix
\input{appendix/a_link}
\input{appendix/b_as}
\input{appendix/c_knapsack}
\input{appendix/d_setup}
\input{appendix/g_interp}
\input{appendix/e_proofs}
\input{appendix/f_evidence}
\input{appendix/f_related}

\end{document}

\section{Conclusion}\label{sec:conclusion}

Viewed through its signed M\"obius expansion, low-bit quantization damage appears to carry irreducible
structure at every order, yet in the Walsh basis roughly 90\% of its variance is order-1, and a
saturating coverage model with one parameter per layer generates both behaviors exactly, the divergent
series and the concentrated spectrum alike. Because the order-$\ge 2$ share equals the best additive
predictor's mean-squared error as a fraction of the variance (\thmone), every result is paired with a
certificate of when additive machinery suffices for prediction, measured on exact lattices or forecast in closed form
from fitted break-rates, and the ordinal defect places the damage within a few percent of a monotone
function of a per-layer sum, so the residual interaction is largely the curvature of the loss scale
rather than structure among layers.

The practical consequences follow. Sensitivities measured in isolation overstate configuration losses
by 71--376\%, by a closed-form inflation factor. The fitted coverage and additive models price a
configuration's loss before deployment, are best or tied against black-box surrogates at small
measurement budgets, and degrade at most $2\times$ under extrapolation to differently quantized
configurations where Gaussian-process surrogates degrade $2$--$5\times$. One Aumann--Shapley path
integral recovers the optimal additive coefficient at a layer-specific effective density, and the
ordering of layers exactly, with no sampled configurations. Used as allocators
at matched memory, the same models
attain the lowest KL among the compared allocators, $17$--$21\%$ below the production ModelOpt baseline
on Qwen3-30B, with consistent margins at 235B and 355B. The
margins
appear where per-unit sensitivities are heterogeneous and budgets are contended, and they vanish in the
served $\ge 4$-bit regime, which is near-lossless by construction. Below four bits, where an RTN
base collapses every allocator, an error-compensating base restores the choice, and allocation decides
whether the served model produces terminating generations at all. 
% The certificate certifies
% prediction, not serving: by \thmone{} it is the exact mean-squared error of the best additive fit
% under the deployment distribution, and the serving limits above are established by the measurements
% themselves.

\section{An interpretable predictor, in and out of distribution}\label{sec:predictor}

Two predictors are fit throughout: the coverage model, and the additive model. Each has $L+1$ parameters, 
fit by least squares in
loss space on $N \approx L$ measured configurations.
The choice of estimator (joint regression rather than per-coefficient sampling), its consistency under
misspecification, and the cost comparison with the ${\sim}100L$ progressive-quantization
evaluations of CoopQ's SPQE estimator and with the fit-free Aumann--Shapley variant are discussed in
App.~\ref{app:setup}.

\subsection{Benchmark}\label{sec:benchmark}

The benchmark task is to predict $f(S)$ for held-out configurations from $N$ measured ones. Per
(model, corpus) cell, every method is fit on $N$ configurations sampled from $\mup$ at $p=0.6$ and
scored by median relative error on 120 held-out configurations, against the surrogate predictors of
Section~\ref{sec:related} (random forest, Gaussian process, RBF kernel model, and CoopQ's
linear-plus-pairwise model), at W4A4 on the six models and two corpora of Table~\ref{tab:master}
(App.~\ref{app:setup}, with the error-versus-$N$ sweep in Fig.~\ref{fig:predn}). At the $N=L/2$ budget,
coverage
or additive is best in 5 of 11 cells and within 0.03 of best in 8, with small margins separating the
leading methods, while isolate-and-sum, at the same measurement cost, fails at 0.38--0.63 error on the
FP4-damaged models. 
% Each cell also reports its certificate, the order-$\ge 2$ share measured as the
% held-out residual of an additive fit (App.~\ref{app:setup}). 

% The three largest shares
% (0.09--0.30) mark the three rows with the worst best-attainable error, and the largest, at the
% crossover value of Section~\ref{sec:dit}, is the one row a black-box surrogate takes by a wide margin.

% With enough in-distribution data the surrogates interpolate best (on 3B-wiki the forest reaches 0.03 at
% $N \approx 8L$ against coverage's 0.08 plateau, while on 0.6B-wiki the Gaussian process leads instead),
% so large-$N$ interpolation selects no single predictor. 
The deciding comparisons are extrapolation to differently quantized configurations,
where coverage degrades at most $2\times$ against the Gaussian process's $2$--$5\times$
(Section~\ref{sec:ood}), allocation, where the surrogates fail outright (Section~\ref{sec:allocmoe}),
and pricing, which only the parametric models provide (App.~\ref{sec:interp}).

% \paragraph{Interpretability enables action.} Unlike the nonparametric surrogates, coverage returns $L$
% per-layer fragility scores $a_i$ and a ceiling $c$. This enables (i) budget pricing: predicting any
% configuration's
% loss before committing, and solving for the cheapest configuration with $\Delta\lce\le\varepsilon$, which
% rankings
% cannot; (ii) targeted precision, applying expensive fixes such as rotation or an extra bit only to
% high-$a_i$ layers; and (iii) mechanistic validation, relating $a_i$ to layer depth, type, and outlier
% kurtosis (Section~\ref{sec:interp}).

\subsection{Out of distribution: stationary kernels do not encode saturation}\label{sec:ood}

On binary configuration vectors the RBF kernel \citep{williams2006gpml} reduces to
$K(S,S') = \exp(-\gamma\,|S\triangle S'|)$, a stationary kernel in the Hamming metric. Under the
uniform measure, kernels of this
form are diagonal in the Walsh basis with eigenvalues that depend only on the order $|T|$ and decay as
$|T|$ grows (\citealp{odonnell2014boolean}; derivation in App.~\ref{app:proofs}), and approximately so
at $p=0.6$, so Gaussian-process regression with this kernel is
shrinkage of high-order Walsh coefficients \citep[Section~2.6]{williams2006gpml}, the same low-order
preference our measurements justify.
This is why the Gaussian process matches our predictors at large $N$ in distribution, where order-1
terms carry most of the damage variance (Section~\ref{sec:walsh}).
The difference emerges out of distribution. 
Stationarity applies the same similarity rule everywhere on
the cube, while saturation does not: under coverage the marginal effect
$f(S\cup\{i\})-f(S)=c\,a_i\prod_{j\in S}(1-a_j)$ shrinks as the context $S$ grows. The stationary prior
does not encode this structure.
% , and its posterior attenuates toward the prior mean exponentially in the
% Hamming distance from the training set.
We test this on the two cells of Table~\ref{tab:seeds} (Qwen3-0.6B and
Qwen2.5-3B, wiki), fitting every method on configurations drawn at $p=0.6$ as before and scoring each
on 120 test configurations drawn at $p=0.85$ 
or at $p=0.35$. At the $N=L/2$ budget, coverage's error grows by
$1.3$--$2.0\times$ and the additive model's by $1.1$--$1.9\times$ relative to their in-distribution
errors across the four (cell, shift) pairs, while the
Gaussian process's grows by $2.1$--$4.9\times$ and the random forest's by $2.0$--$3.9\times$
(Fig.~\ref{fig:ood}). 

% The same protocol at W4A4 was also measured and does not separate the methods,
% since those cells are nearly saturated: the $p=0.85$ shift is easy for every method, and at $p=0.35$
% on the 3B cell every method's absolute error becomes large (coverage smallest at 0.71, the Gaussian
% process at 19.4). A shift test separates methods only between those two poles, where the shifted losses change enough
% to be nontrivial but not so much that every method fails. The W16A4 cells sit in that regime, and the
% extrapolation claims of this section are therefore scoped to them.
% Coverage's worst error under either shift (0.22) is below the Gaussian process's
% best (0.26). 
The heavily quantized direction is the one allocation exercises, hence an explanation for its
performance in allocation.
% The greedy of
% Section~\ref{sec:allocation} prices configurations more heavily quantized than any it was fit on, so a
% predictor's error under that shift sets allocation quality, and
% Section~\ref{sec:allocmoe} measures the failure directly.

\begin{figure}[t]
  \centering
  \includegraphics[width=\linewidth]{fig4_ood.png}
  \caption{Fit at $p=0.6$, tested at $p=0.85$ and $p=0.35$. Coverage degrades $1.3$--$2.0\times$ and the additive fit $1.1$--$1.9\times$ relative to
  their in-distribution errors, while the
  Gaussian process degrades $2.1$--$4.9\times$ and the random forest $2.0$--$3.9\times$; the
  stationary priors of the surrogates do not encode saturation.}
  \label{fig:ood}
\end{figure}

\subsection{A diffusion transformer}\label{sec:dit}

The benchmark repeats on the Qwen-Image-2512 diffusion transformer ($L=60$ blocks, per-block NVFP4),
with $f(S)$ the mean squared error between the quantized and full-precision model outputs at the first
denoising step (per-budget results in App.~\ref{app:setup}, Table~\ref{tab:dit}). The measured
order-$\ge 2$ share is 0.36, above every LLM cell (0.016--0.30, Tables~\ref{tab:master}
and~\ref{tab:seeds}). At small budgets, coverage and random forest edge the other methods,
with performance converging across methods as budget size increases.

% At small budgets coverage attains the lowest error of every method except the
% random forest: 0.15 at $N=L/2$ against the Gaussian process's 0.37, with additive at $3.2\times$
% coverage's error at $N=L$, and both matching coverage only from $N \approx 4L$. The forest, in
% contrast, matches coverage immediately (0.142 against 0.151 at $N=L/2$), and at $N \ge 4L$ additive,
% coverage, the forest, and the Gaussian process converge within ${\approx}0.02$ of each other.

% \paragraph{The crossover is monotone in the certificate (Fig.~\ref{fig:crossover}).} Across cells, the
% (RF $-$ coverage) error gap falls as the share rises: $+0.048$ at share 0.05 and $+0.011$ at 0.26 (the
% W16A4 cells of Table~\ref{tab:seeds}, $N=L$), $-0.009$ at 0.30 (the 30B cell of
% Table~\ref{tab:master}, $N=L/2$), and $-0.009$ at 0.36 (the DiT, $N=L/2$), crossing zero at a share of
% roughly 0.3. The forest can fit the higher orders but pays for them in data, while coverage's $L{+}1$
% parameters are the matched model when per-layer terms dominate. By \thmone{}, no predictor can improve
% on the best additive one by more than the order-$\ge 2$ energy, so a small share leaves the forest
% nothing to gain and a large share permits the gain it realizes. The certificate is computable from the
% same measured configurations, so the choice of model class can be made before committing to one.
%
% \begin{figure}[t]
%   \centering
%   \includegraphics[width=0.7\linewidth]{fig10_cov_rf_crossover.png}
%   \caption{The certificate predicts the parametric-to-nonparametric crossover: the (RF $-$ coverage)
%   small-budget relative-error gap against each cell's order-$\ge 2$ share, for the four cells and
%   budgets listed in the text. The gap falls monotonically and crosses zero at a share of
%   $\approx 0.3$. By \thmone{}, improvement over the best additive predictor is capped by the share.}
%   \label{fig:crossover}
% \end{figure}

\section{Future work}\label{sec:future}

\paragraph{Data randomness in the theory.} We treat $f(S)=\Delta L$ as deterministic, but it is an
expectation over the calibration distribution, with data-dependent break-rates. A two-source variance
decomposition, layer-selection variance (the $p$-biased Walsh account) plus data-sampling variance, would
give a principled law connecting calibration size to prediction variance and a data-aware certificate.

\paragraph{Hardware- and framework-aware allocation.} Effective
bits equal memory is the wrong cost for realized speedup, which is a property of the serving stack: fusion
granularity sets the alphabet (one format per fused qkv or expert block), 
kernel availability determines if certain bit-widths (e.g. sub-4) are servable.
Co-designing the cost function of the allocator with the serving stack would lead to better realized speed gains.

% , and the batch and parallelism regime
% can flip the sign of the speedup (communication-bound decode buys nothing, GEMM-bound prefill helps).
% Replacing effective bits with a measured per-(family, format, fusion-group, batch, parallelism)
% latency-delta cost table, with fusion and kernels as hard constraints, would co-design the allocator with
% the serving stack and bridge predicted quality to realized speedup; under such constraints the clean
% per-bit knapsack of App.~\ref{app:knapsack} becomes a constrained program, and global methods (ILP,
% relaxation-plus-rounding) replace the greedy.

% \paragraph{Extreme quantization.} 
% % The GPTQ ladder delivers a served sub-4 result on one model and one
% % task (Section~\ref{sec:allocsub4}). The open items are replication of the non-termination onset and of
% % the allocator ordering on further models and tasks, a base quantizer that survives
% % 2 bits
% % (QuIP\#-class incoherence processing, since at 2 bits scale granularity dominated error compensation in
% % our comparison), and finer grids between 2 and 3 bits where the rationing question is sharpest.
% Composing our allocation algorithm with base quantizers which mitigate extreme quantization (e.g. GPTQ, QuIP\#) is an 
% open application.

\paragraph{Composing allocation with outlier and rotation methods.} Allocation
chooses how many bits each unit gets, while rotations (QuaRot, SpinQuant, QuIP\#), SmoothQuant, and
llm.int8() change how much a unit is damaged at a given bit-width, 
% lowering the per-unit break-rate and,
% under a shared scale, the cross-layer coupling itself (SmoothQuant moves the order-$\ge 2$ share from
% 0.67 to 0.46, and outlier clipping from 0.67 to 0.16, Section~\ref{sec:scope}). 
Our framework operates on
whatever base quantizer produces $f(S)$, so the two compose: allocate over a rotated or outlier-handled
base. 
% The open question is the sign of the combination, which the theory makes precise. Outlier handling
% that flattens heterogeneity, as a good rotation is designed to, would shrink the headroom allocation
% exploits, whereas outlier handling that lowers a uniform floor
% leaves the relative heterogeneity intact. The certificate and the spread of the
% break-rates quantify this directly: measuring them before and after the orthogonal method predicts whether
% composition compounds the gains or cannibalizes them, and targets the expensive fix at high-$a_i$ units
% (Section~\ref{sec:scope}). This converts allocation versus outlier methods into allocation over outlier
% methods, with a predictive rule for when it pays.

\section{The measured structure of quantization damage}\label{sec:walsh}

This section reports the two measurements on which the rest of the paper builds: the decomposition of the
damage variance by order, and the ordinal defect, which measures how much of that structure survives every
increasing transformation of the loss.

\paragraph{Set-up.} For a contiguous block of $L$ layers (finer units are quantized in
Section~\ref{sec:allocation}), centered at the network's mid-depth
(App.~\ref{app:setup}), we measure $f(S)=\Delta\lce$ for the full subset lattice of $2^L$ subsets
$S$, running the model with exactly the layers in $S$ quantized. The quantizer rounds stochastically, so
each $f(S)$ is the average of several independent rounding draws (16 for the lattices in this section;
App.~\ref{app:setup}). A measure over configurations must be fixed before variances or prediction
errors can be defined, and it should weigh the configurations that matter in practice, which are mostly
quantized. We therefore analyze $f$ under the \emph{deployment distribution} $\mup$, the product measure
in which each layer is quantized independently with probability $p$, and take $p=0.6$ throughout unless
stated otherwise. We expand the measured $f$ in the
\emph{Walsh basis} $\{\chi_T\}$, with coefficients $\hat f(T)$ \citep{odonnell2014boolean}. We denote the
\emph{order-$k$ energy} $W_k = \sum_{|T|=k} \fhat(T)^2$, and by Parseval's identity, $\Var(f)=\sum_{k\ge 1}W_k$, so the ratio $W_k/\Var(f)$ is the
fraction of the loss variance attributable to order-$k$ structure. The following identity gives these energies their
operational meaning.

\begin{customthm}{1}[additive error equals order-$\ge 2$ energy]\label{thm:one}
For any $f$ on $\{0,1\}^L$ and any $p\in(0,1)$, the best degree-$\le 1$ 
predictor $f_1$ under $\mup$ has squared error exactly equal to the high-order Walsh energy:
$\E_{\mup}\bigl[(f-f_1)^2\bigr] = \sum_{k\ge 2} W_k$.
\end{customthm}

\begin{proof}
The degree-$\le 1$ functions are exactly $\mathrm{span}\{\chi_\emptyset, \chi_{\{1\}}, \ldots,
\chi_{\{L\}}\}$, least squares in $L^2(\mup)$ is the orthogonal projection onto that span, and by
Parseval the squared residual is $\sum_{|T|\ge 2}\fhat(T)^2$.
\end{proof}

This is the $d=1$ case of low-degree learning \citep{odonnell2014boolean,lmn1993,sobol2001}.

\paragraph{Result (Fig.~\ref{fig:walsh-energy}).} Roughly 90\% of the variance is order-1, with
order-$\ge 2$ contributing 7--15\% across models and datasets. On a 10-layer Qwen3-8B block,
order-1 accounts for $87.3\%$ of the variance as measured and $98.1\%$ once the noise floor
contributed by the stochastic-rounding draws is subtracted (the split-half estimate of
App.~\ref{app:ordinal}), and on an 8-layer Qwen3-0.6B block for
$93.8\%$ and $99.2\%$. Across $\{$Qwen3-0.6B, Qwen2.5-3B, Qwen3-4B, Qwen3-8B, Llama-3.2-3B,
Qwen3-30B-A3B$\}\times\{$wiki, code$\}$ (WikiText-2 and a code corpus, App.~\ref{app:setup}) the
sampled order-1 share is 85--93\% when whole layers are quantized at FP4. The concentration is not an
artifact of whole-layer units or of small scale (Fig.~\ref{fig:walsh-energy}): on exact $2^8$ lattices
over eight consecutive linear units at W4A4, the order-1 share is 97--99\% on models from 0.6B to
Qwen2.5-72B, and 90.7\% on gpt-oss-120B.

\paragraph{The intercept diagnostic (Fig.~\ref{fig:walsh-intercept}).} An additive model with no
intercept, which is the constraint that aggregating isolated sensitivities imposes ($f(\emptyset)=0$),
reaches $R^2=0.13$ on the 0.6B block, where FP4 damage is large, against $0.78$ on the 8B block, where it
is small. Relaxing this constraint raises both to $0.94$ and $0.87$. Both fits are least squares under
$\mup$ on the full lattice.

The remaining questions are why the small order-$\ge 2$ content is structured the way it is and what
generates it. The answer we propose is saturation (Section~\ref{sec:theory}). Before turning to it, we ask how much of the
order-$\ge 2$ energy reflects structure among layers at all.

\subsection{The ordinal defect: additivity up to an increasing transformation}\label{sec:ordinal}

The order-$\ge 2$ energy has two caveats as an indicator of structure among layers. First, it is
measure-relative. 
Second, and more fundamentally, it is not
invariant under a change of coordinates on the loss values, since replacing $f$ by $\psi \circ f$ for a
strictly increasing $\psi:\mathbb{R}\to\mathbb{R}$ can change the energy if $\psi$ is curved, 
so an exactly additive $f$ could acquire nonzero higher-order energy. A principled
notion of irreducible structure should therefore quantify over all strictly increasing $\psi$. For
binary factors this is possible, because additivity up to an increasing transformation is a purely ordinal
property.

\begin{customthm}{2}[additivity up to an increasing transformation is ordinal]\label{thm:ordinal}
$f(x) = g(b\cdot x)$ for some $b\in\mathbb{R}^L$ and strictly increasing $g$ if and only if some additive
index reproduces $f$'s ranking of subsets, that is, $f(x) < f(y) \iff b\cdot x < b\cdot y$ for all
$x,y \in \{0,1\}^L$.
\end{customthm}

The theorem is the finite binary case of additive conjoint measurement
\citep{lucetukey1964conjoint,debreu1960topological}. App.~\ref{app:link} gives the
short proof. Throughout, an \emph{additive index} is a linear score $x \mapsto b\cdot x$ used only
through the ordering it induces on configurations. Whether such an index exists is a linear program, with one inequality
$b\cdot(x-y)>0$ per
ordered pair \citep{scott1964measurement}. The smallest witnesses of infeasibility are \emph{rank flips},
pairs of comparisons that order the same two layers oppositely in two contexts:
$f(S\cup i) > f(S\cup j)$ but $f(T\cup i) < f(T\cup j)$ forces both $b_i>b_j$ and $b_i<b_j$, so a single
true flip refutes every strictly increasing $\psi$. We call the minimal fraction of ranking
constraints that any additive index must violate the \emph{ordinal defect} of $f$.

\paragraph{Measurement.} On full $2^8$ configuration lattices, measured with 16 stochastic-rounding
draws per configuration (App.~\ref{app:setup}), we bound the defect from both sides, certifying each
statistic against measurement noise. For the lower bound we count rank flips. Splitting the 16 draws
into two halves gives two independent estimates of every configuration's loss, and a flip is counted
only if it is detected above the noise floor on the first half and recurs with the same sign on the second. The surviving
count is then read against the count that the identical pipeline produces on a matched-noise null, a
synthetic lattice constructed to contain no true flips (the coverage model of Section~\ref{sec:theory} fitted
to the measured losses, plus noise at each configuration's measured level). For the upper bound we
exhibit an index, since the fraction of ranking
constraints violated by any single index bounds the defect from above. The exhibited index minimizes the total
slack of the constraints of one half of the draws, again a linear program, and its violated fraction is
evaluated both on that half and on the held-out half's constraints, with the larger of the two quoted
(App.~\ref{app:ordinal}, per-model results in
Table~\ref{tab:ordinal}).
In the deployed per-block regime the certified ordinal defect is
at most $1.4\%$ (the largest cell, Llama-3.2-3B). The energy counterpart agrees: minimizing the order-$\ge 2$
share over all strictly increasing transformations of the loss (the estimator $E^*$ of
App.~\ref{app:ordinal}) yields an upper bound of $0$--$5\%$
across the deployed cells (last column of Table~\ref{tab:ordinal}). The zeros are not blindness of the procedure. Applied to the non-deployed
per-tensor regime of Section~\ref{sec:scope}, the same procedure finds flips well in excess of its
matched-noise null (74 against 18, Table~\ref{tab:ordinal}), so genuine reversals are detected where a
known mechanism produces them, and even there an additive
index orders at least $98\%$ of the comparisons above the noise floor correctly. 
The measured
$f$ is therefore, to that accuracy, of the form $g\bigl(\sum_{i\in S} b_i\bigr)$ for a strictly increasing $g$, and
the apparent high-order energy reflects the curvature of $g$ rather than structure among layers.



\begin{figure}[t]
  \centering
  \begin{subfigure}[t]{0.615\linewidth}
    \centering
    \includegraphics[width=\linewidth]{fig1a_order1_share.png}
    \caption{Order-1 share, every exact lattice.}
    \label{fig:walsh-energy}
  \end{subfigure}\hfill
  \begin{subfigure}[t]{0.365\linewidth}
    \centering
    \includegraphics[width=\linewidth]{fig1b_intercept.png}
    \caption{The intercept diagnostic (in sample).}
    \label{fig:walsh-intercept}
  \end{subfigure}
  \caption{(\subref{fig:walsh-energy}) Order-1 share of the damage variance on every exact lattice
  measured, under the deployment measure $\mup$, $p=0.6$. Red: the two whole-layer FP4 lattices of this
  section (Qwen3-0.6B, 8 layers, $2^8$ configurations; Qwen3-8B, 10 layers, $2^{10}$), noise-corrected
  as in App.~\ref{app:ordinal}. Teal: exact $2^8$ lattices at per-linear W4A4 granularity (NVFP4,
  block 16), each over eight consecutive mid-depth linear units, with damage measured as the KL
  divergence to the full-precision model on a WikiText-2 calibration set (light bars: code corpus).
  The share is 96.9--99.7\% on dense models from 0.6B to 72B and 90.7\% on the sparse-MoE
  gpt-oss-120B, where fused expert blocks add coupling.
  (\subref{fig:walsh-intercept}) The additive fit without an intercept, the constraint that
  isolated-sensitivity aggregation imposes ($f(\emptyset)=0$), against the same model class with
  a free intercept, which by \thmone{} is the degree-$\le 1$ projection; least squares under $\mup$ on
  the two whole-layer lattices.}
  \label{fig:walsh}
\end{figure}

\section{Why isolated sensitivity mis-prices a configuration: the wrong anchor}\label{sec:anchor}

Throughout, a predictor's \emph{anchor} is the
configuration, or distribution over configurations, at which its coefficients are measured: isolated sensitivities
are anchored at the unquantized network,
while the coefficients of a fitted additive model are anchored at the deployment distribution.

\paragraph{The measurement.} An estimate of a configuration's loss built from per-layer sensitivities
alone is the sum $\sum_{i\in S} f(\{i\})$, which we call \emph{isolate-and-sum} throughout. 
 On the fully quantized configuration of the
exact whole-layer FP4 lattices (the two lattices of Section~\ref{sec:walsh} and the four
certification cells of Table~\ref{tab:ordinal}) this aggregation overstates the measured loss by
71--346\%. The mis-pricing persists on full networks. Predicting $f$ on
configurations that protect a sampled subset of layers and quantize the rest (all 36 Qwen3-8B and 48
Qwen3-30B-A3B layers, protection sizes 1 to
$L/2$), isolate-and-sum's median relative error is
51\% on Qwen3-8B and 203\% on Qwen3-30B-A3B, while additive and coverage models fitted on 80
sampled configurations attain 2--3\% on the same held-out protection sets. HAWQ-v2, scoring each layer by its Hessian trace (Hutchinson estimator) times its mean squared
quantization perturbation \citep{dong2020hawqv2}, inherits the same anchor: summed over each of the
two Section~\ref{sec:walsh} lattices it overstates the fully quantized corner by 89\% (Qwen3-8B) and
376\% (Qwen3-0.6B), against isolate-and-sum's 155\% and 250\% on the same corners.

\paragraph{Marginal scale.} The failure is caused by two distinct order-1 coefficients.
The \emph{in-context} slope, the coefficient of the optimal additive fit, is
$w_i = \E_p[f\mid x_i=1] - \E_p[f\mid x_i=0]$, the marginal effect of layer $i$ averaged over deployment
configurations. The \emph{isolated} slope, the quantity that HAWQ-style scores approximate,
is $s_i = f(\{i\})$, the marginal at the
unquantized network. Under coverage, $s_i = c\,a_i$ while $w_i = c\,a_i\prod_{j\ne i}(1-p a_j)$, so
\[
  \frac{s_i}{w_i} = \frac{1}{\prod_{j\ne i}(1-p a_j)} \approx e^{\,p\sum_j a_j},
\]
a ratio that grows exponentially with $p\sum_j a_j$. We validate the formula directly on the
full $2^8$ lattices, comparing the measured ratio $s_i/w_i$ 
against the value the fitted coverage model predicts. The measured medians are 1.9 / 3.0 / 3.9 / 9.2
(Llama-3.2-3B / Qwen3-30B-A3B / Qwen3-8B / Qwen3-0.6B) against predicted
1.8 / 2.3 / 2.5 / 7.7, with more detail in
App.~\ref{app:knapsack}. 
The two-slope identity is a one-line corollary of Theorems~1 and~4 (App.~\ref{app:proofs}).
% A small order-$\ge 2$ share bounds the mean-squared error of the fitted additive model (\thmone{}), but
% it says nothing about the isolated-anchor sum, whose coefficients are inflated by the factor above
% regardless.

% Section~\ref{sec:allocation} uses this mis-pricing directly: allocators scored with the in-context
% coefficient outperform allocators scored at the unquantized-network anchor, at matched memory. 

\section{Allocation at matched memory}\label{sec:allocation}

\subsection{Setup}\label{sec:allocsetup}

Given a memory budget in effective bits,
$\mathrm{eff} = \sum \mathrm{numel}\cdot\mathrm{bits} / \sum \mathrm{numel}$, the allocator assigns each
unit (a weight tensor, or a fused group that must share one serving format, App.~\ref{app:setup}) a
bit-width in $\{16,8,4\}$ to minimize predicted $\Delta\lce$ (Section~\ref{sec:allocsub4} extends the
alphabet below four bits). Under the coverage model, predicted damage increases
in the summed log-headroom costs $b_i$ of the demotions taken, so minimizing it under a bit budget is a
knapsack over those per-demotion costs (App.~\ref{app:knapsack}). 

% Demotion prices are estimated in deployment context, from a
% coverage or additive fit on sampled configurations or fit-free from one
% Aumann--Shapley path integral per demotion level, written AS-cov and AS-add (App.~\ref{app:as}).

\paragraph{Converting the baselines into allocators.} The
Section~\ref{sec:predictor} baselines are predictors, so each is fit on the same $N=300$
sampled configurations and run through the same knapsack over its predicted marginals (\code{fit-add} for
the fitted additive model, \code{gp-alloc} and \code{rf-alloc} for the surrogates). 
% Two incumbents complete the set: CoopQ/SPQE
% \citep{impq2025}, and MC-MoE
% \citep{mcmoe2025}, which ranks experts by a routing-frequency-times-mass prior and uses no damage
% measurements. 
We fit the coverage and additive scores either with least squares or an Aumann-Shapley path integral.
The Aumann--Shapley score is the loss gradient averaged over increasingly quantized contexts, and under the
coverage model it equals the optimal in-context coefficient evaluated at an effective deployment
density (App.~\ref{app:as}). 
% The metric is the KL
% divergence from the bf16 model's next-token distribution to the quantized model.


\subsection{Serving against vendor NVFP4 checkpoints}\label{sec:alloc27b}

On the dense Qwen3.6-27B we compare the additive allocation at 5.41 effective bits, derived from
measured per-unit step damages on a task-aligned calibration set (App.~\ref{app:setup}), against the
two public NVFP4 checkpoints, NVIDIA's at 6.31 effective bits and Unsloth's at 6.73, all served by
vLLM on one B200 at concurrency 128. Ours and Unsloth's checkpoints serve W4A4 NVFP4, quantizing
activations at runtime against per-block scales; NVIDIA's published checkpoint keeps its NVFP4
linears weight-only (W4A16). The additive configuration matches the vendor checkpoints
downstream (Table~\ref{tab:race27b}), and the smaller footprint is
realized as speed: it is the fastest system at every prompt/generation shape, $1.5$--$2.0\times$ the
bf16 throughput and 8--24\% above the Unsloth checkpoint (Fig.~\ref{fig:race27b}).

\begin{table}[t]
\centering\small
\caption{Qwen3.6-27B served accuracy (pass@1; AIME25 averaged over 8 seeds; thinking sampling).}
\label{tab:race27b}
\begin{tabular}{lcccc}
\toprule
 & eff.\ bits & AIME25 & GPQA & MMLU-Pro \\
\midrule
bf16 & 16 & 90.0 & 83.8 & 86.3 \\
NVIDIA NVFP4 & 6.31 & 90.8 & 67.2 & 85.8 \\
Unsloth NVFP4 & 6.73 & 90.4 & 82.3 & 86.7 \\
additive & 5.41 & 90.0 & 79.8 & 86.5 \\
\bottomrule
\end{tabular}
\end{table}

\begin{figure}[t]
  \centering
  \includegraphics[width=\linewidth]{fig_race27b.png}
  \caption{Qwen3.6-27B output throughput at four prompt/generation shapes,
   with the speedup over bf16 marked on the additive bars. The additive
  configuration is fastest at every shape at the lowest effective bits.}
  \label{fig:race27b}
\end{figure}

The comparison repeats at 753B. On GLM-5.2 we compare the additive allocation (4.84 stored weight
bits per parameter, 474~GB) against NVIDIA's published NVFP4 checkpoint (4.76 bits, 465~GB) and the
FP8 base model, served by SGLang on $4\times$B200 at concurrency 8. Both NVFP4 configurations are W4A4.
 Accuracy matches at matched
stored size (Fig.~\ref{fig:glm52}), and the additive configuration reads 28.3~GB of weights per
decoded token against the checkpoint's 45.9, so it decodes $1.19\times$ faster: the allocation
places its low-precision bits on the units decode reads most often.

\begin{figure}[t]
  \centering
  \includegraphics[width=0.5\linewidth]{fig_glm52.png}
  \caption{GLM-5.2 (753B) decode throughput, SGLang on $4\times$B200 at concurrency 8, with total
  stored weight and weight bytes read per decoded token (batch size 1) below each bar. Accuracy is
  matched: MMLU 0.842 (additive) against 0.838 (NVIDIA NVFP4) and 0.864 (the FP8 base), GPQA diamond
  0.470 against 0.455 and 0.465 (zero-shot letter-logit scoring; 228-question stratified MMLU subset,
  198 GPQA questions).}
  \label{fig:glm52}
\end{figure}
\subsection{MoE language models}\label{sec:allocmoe}

We benchmark our methods against the incumbent methods. The certificate accompanying these results
depends on the units it is measured over. On the eight whole-layer FP4 block lattices of Qwen3-30B
(Section~\ref{sec:tauval}) the measured order-$\ge 2$ share has median 0.013. At the allocation
granularity itself (per-linear units over the $\{16,8,4\}$ alphabet) no lattice is enumerable, and
the split-half cross-estimate (coefficient products of two disjoint halves, as in
App.~\ref{app:ordinal}) from the 300 sampled configurations of the fitted baselines, unbiased
but not constrained to be nonnegative, gives $-0.09 \pm 0.10$ on the 30B pool and $0.24 \pm 0.14$ on
the 235B pool. At $N=300$ the sampling error is of the order of the share itself, so the
allocation-granularity certificate is an order-of-magnitude statement, not a precise value. 
% The
% $\tau$-forecasts from fitted break-rates are
% ${\sim}10^{-4}$ at 235B and GLM-4.6, and where a lattice exists to check this pipeline it understates
% the measured share (App.~\ref{app:setup}). 
Per-method values are in
Tables~\ref{tab:30bb} and~\ref{tab:235c}, each with a paired test over the 16 shared evaluation
sequences, Holm-corrected within each budget row at the 5\% level.
%  the
% budget-by-budget curves in Fig.~\ref{fig:pareto}. 
On Qwen3-30B the additive- and coverage-scored allocators lead at every measured budget.

% Only the 235B eff-8 corner, near-lossless with nothing to allocate, goes to all-FP8, while the more
% damage-sensitive 30B differentiates even there (Section~\ref{sec:allocwhere}: the contended range is
% set by the damage level, not by effective bits alone).

% \paragraph{Reading the KL margins.} The servable $\ge 4$-bit regime is near-lossless downstream
% (Section~\ref{sec:scope}), so these matched-memory KL margins are a leading indicator rather than a
% realized accuracy gap: a distributional safety margin for settings that are not near-lossless (harder
% tasks, tighter formats, longer horizons). The margins are consistent across three MoE scales (Qwen3-30B
% and 235B, GLM-4.6 at 355B) and every contended budget. Below four bits the margin is realized, deciding whether
% the served model functions at all (Section~\ref{sec:allocsub4}).



\begin{table}[t]
\centering
\caption{Qwen3-30B allocation comparison at matched effective bits, calibration KL on the task-aligned
corpus of App.~\ref{app:setup} (BFCL and RULER prompts).
$*$ marks methods
significantly different from AS-additive under a paired $t$-test over the 16 shared sequences,
Holm-corrected within each budget row at the 5\% level. $^\dagger$ marks MC-MoE's per-expert form, which assigns each
expert its own format and cannot be loaded by the fused serving kernels, so the fusion column is its
servable form. int8-outlier ranks units by an llm.int8-style activation-outlier severity score.
Marginal 95\% intervals on the entries span
$\pm.001$--$.007$.}
\label{tab:30bb}
\resizebox{\textwidth}{!}{%
\begin{tabular}{ccccccccccccc}
\toprule
eff & naive-unif & naive-routed & ModelOpt & int8-outlier & CoopQ (1{,}205 ev) & MC-MoE per-exp$^\dagger$ & MC-MoE fusion &
fit-add (300 ev) & gp-alloc (300 ev) & rf-alloc (300 ev) & AS-cov & AS-add \\
\midrule
4.2 & .0480$^*$ & .0367 & .0429$^*$ & .0428$^*$ & .0419$^*$ & .0359 & .0359 & .0372 & .0382$^*$ & .0538$^*$ & .0365 & \textbf{.0353} \\
4.5 & .0344 & .0344 & .0383$^*$ & .0327 & .0393$^*$ & .0318 & .0337 & \textbf{.0310} & .0319 & .0486$^*$ & .0317 & .0321 \\
5.0 & .0335$^*$ & .0344$^*$ & .0331$^*$ & .0310$^*$ & .0347$^*$ & .0275 & .0311 & .0284$^*$ & .0295$^*$ & .0443$^*$ & .0266 & \textbf{.0264} \\
6.0 & .0318$^*$ & .0344$^*$ & .0253$^*$ & .0253$^*$ & .0308$^*$ & .0232 & .0255 & .0229$^*$ & .0235$^*$ & .0421$^*$ & .0202 & \textbf{.0200} \\
8.0 & .0137$^*$ & .0344$^*$ & .0215$^*$ & .0177$^*$ & .0137$^*$ & .0192 & .0119 & .0185$^*$ & .0201$^*$ & .0371$^*$ & \textbf{.0106} & .0109 \\
\bottomrule
\end{tabular}}
\end{table}

The predictor skill of a black-box surrogate does not transfer: \code{gp-alloc} trails the AS and
fitted allocators by 3--18\% at the contended budgets (the extrapolation failure of
Section~\ref{sec:ood} in allocation form), and \code{rf-alloc} falls below naive-uniform.
CoopQ trails AS by $+19$ to $+54\%$
despite 1{,}205 sampled KL evaluations against AS's zero, and it scores whole layers rather than the
240 per-linear units of the other methods, the only granularity its measurement design affords, since
its permutation-walk cost scales with the unit count.
% , since
% walk cost and its pairwise surrogate scale with the unit count (${\approx}24{,}000$ evaluations at its
% published budget, App.~\ref{app:setup}).
% Its
% pairwise term $K$ contributes nothing: allocating from its linear term alone changes the KL by at most
% .0003 at any contended budget, which the certificate (order-$\ge 2$ share .013) indicated in advance.

% Serving the same eleven methods' configurations across these budgets (allocations derived on the
% evaluation-task corpus of Table~\ref{tab:sub4kl}, weight-only RTN checkpoints, five benchmarks)
% realizes none of these margins: 
% Serving the same methods' configurations at these budgets realizes none of these KL margins:
Table~\ref{tab:ge4served30} shows AIME-2025 scores for the evaluated methods on Qwen3-30B,
with every method between 66 and 76 against the bf16 anchor 73.3 and no significant ordering at any
budget. This finding survives both W4A4 and W4A16 quantization regimes.
% The finding survives the serving channel: re-serving the additive, coverage, and ModelOpt
% configurations at effective bits 4.5 and 5.0 as full W4A4 ModelOpt exports (activations quantized at
% runtime, served by TensorRT-LLM; q/k/v tied to the group maximum for fused serving, shifting the AS
% budgets by at most $+0.09$ effective bits) gives 68.8--72.9 across the six cells, again with no
% significant method ordering at either budget.


% (d) MC-MoE's routing-frequency-times-mass prior involves no damage measurement. It nearly matches AS at
% eff $\le 4.5$ ($+0.3$ to $+1.7\%$, in its unservable per-expert form) but trails by $+4\%$ at eff 5.0
% and $+16\%$ at eff 6.0.
% (e) Per-expert MC-MoE's KL is $0/5.6/11.6/9.0\%$ lower than its fusion-constrained form's across the
% contended budgets, the measured price of fused serving. Per-expert kernels would recover roughly 10\%
% and still trail AS at eff $\ge 5$.
% (f) int8-outlier is significantly worse than AS-additive at four of the five budgets ($+17$ to
% $+62\%$, paired test), indistinguishable only at eff 4.5: an activation-outlier score is an isolated
% magnitude proxy, the Section~\ref{sec:anchor} mis-pricing in allocation form.
% (g) At the loose eff-8 corner, the binary $\{8,4\}$ alphabets degenerate to all-FP8 (.0137), ModelOpt
% spends the same memory for .0215, and the three-level allocators promote a few sensitive units to 16
% bits and improve on the degenerate point (AS-cov .0106, MC-MoE fusion .0119), while naive-routed, whose
% rule fixes routed experts at the lowest level, stays flat at .0344 regardless of budget.

\paragraph{Scale replication (Qwen3-235B, Table~\ref{tab:235c}).} The comparison was re-run at
$8\times$ scale on Qwen3-235B. The in-context family again carries the contended budgets: AS-additive
is best at eff 4.2, the fitted additive and coverage variants at 4.5--6.0, and every method outside
the family is significantly worse at eff 4.2, 4.5, and 6.0 (at 5.0, naive-uniform, ModelOpt, and
\code{gp-alloc} are not separated after the Holm correction). 
% The loose eff-8 corner goes to
% naive-uniform, where near-lossless budgets leave little to allocate. 
The comparison extends once more
to
GLM-4.6 (355B): a three-level $\{16,8,4\}$ allocation over its 43{,}364 units gives AS-coverage
calibration KLs of 0.059, 0.045, and 0.036 at effective bits 5, 6, and 7, against ModelOpt's 0.088,
0.071, and 0.045, margins of 34, 37, and 20\%, with AS-additive within 0.006 of AS-coverage at
every budget. 

\begin{table}[t]
\centering
\caption{Qwen3-235B, all methods measured jointly on the task-aligned corpus of
App.~\ref{app:setup}, with paired significance vs.\ AS-additive over the 16 shared sequences
($*$ = significant after Holm correction within each budget row at the 5\% level).}
\label{tab:235c}
\resizebox{\textwidth}{!}{%
\begin{tabular}{ccccccccccc}
\toprule
eff & naive-unif & naive-routed & ModelOpt & CoopQ (942 ev) & fit-add & fit-cov & gp-alloc & rf-alloc &
AS-cov & AS-add \\
\midrule
4.2 & 0.0196$^*$ & 0.0171$^*$ & 0.0222$^*$ & 0.0205$^*$ & 0.0186$^*$ & 0.0183$^*$ & 0.0185$^*$ & 0.0232$^*$ & 0.0172$^*$ & \textbf{0.0163} \\
4.5 & 0.0161$^*$ & 0.0163$^*$ & 0.0180$^*$ & 0.0190$^*$ & \textbf{0.0149} & 0.0158 & 0.0163$^*$ & 0.0211$^*$ & 0.0161$^*$ & 0.0151 \\
5.0 & 0.0144 & 0.0163$^*$ & 0.0151 & 0.0185$^*$ & \textbf{0.0129} & 0.0133 & 0.0141 & 0.0198$^*$ & 0.0132 & 0.0139 \\
6.0 & 0.0141$^*$ & 0.0163$^*$ & 0.0128$^*$ & 0.0157$^*$ & 0.0102 & \textbf{0.0094}$^*$ & 0.0116$^*$ & 0.0188$^*$ & 0.0112$^*$ & 0.0101 \\
8.0 & \textbf{0.0068}$^*$ & 0.0163$^*$ & 0.0096 & 0.0068$^*$ & 0.0083$^*$ & 0.0083 & 0.0099 & 0.0154$^*$ & 0.0090 & 0.0091 \\
\bottomrule
\end{tabular}}
\end{table}



% Every structural finding replicates. The in-context family leads every contended budget with members
% within $\pm.001$--$.003$ of each other, AS-additive is uniquely best at the tightest budget, 
% and every baseline outside the family is significantly worse at eff 4.2, 4.5,
% and 6.0 (at 5.0, naive-uniform and \code{gp-alloc} fall within noise). 
% CoopQ's deficit grows with
% scale, from $+19$ to $+54\%$ at 30B to $+26$ to $+55\%$ here, at 942 evaluations against zero.
% The deficit follows from its measurement design rather than its surrogate class: our models fit on the
% 942 configurations its permutation walks visited reproduce it, while the same models fit on 300
% independently sampled configurations land within $.002$ of AS (not shown).

% The eff-8 budget again belongs to
% near-uniform configurations, as scoped in Section~\ref{sec:scope}. Within the family, significance
% resolves the near-ties: at 30B eff 5.0/6.0 the fit-free AS is significantly better than its
% 300-evaluation fitted siblings, while at 235B eff 5.0/6.0 fitted-coverage leads by a starred
% $.0006$--$.0007$. The spread within the family is at noise scale, and the gap between the family and
% the rest is not.



\subsection{Allocation below four bits}\label{sec:allocsub4}

% Below four bits allocation stops being a proxy-metric comparison and starts deciding whether the served
% model functions at all. The regime is productive: the allocated configurations serve at bf16-level
% gsm8k accuracy at effective bits 3.5 and continue solving the task down to 2.75, where the
% gradient-scored configuration at the same memory no longer produces terminating generations
% (Table~\ref{tab:sub4served}). 
Below four bits the base quantizer decides whether allocation has anything to work with: under
round-to-nearest quantization, every allocator we measured collapses at every sub-4 budget regardless of model
scale, so we use an error-compensating base. We construct a
per-unit GPTQ \citep{frantar2022gptq} ladder on Qwen3-30B, quantizing every unit at 4, 3, and 2 bits
(weight-only; details in App.~\ref{app:setup}). 

% Error compensation opens the 3-bit rung: with all units at
% one level, the calibration KL is 0.054 at GPTQ-4, 0.202 at GPTQ-3, and 13.6 at GPTQ-2, so the usable
% ladder is $\{4,3,2\}$ with 2 bits as a level to be rationed. The 13.6 is carried by the attention
% projections, which hold only 0.8\% of the quantized weight: with those kept at 4 bits, all-2 over the
% expert blocks reaches KL 1.29, the floor of Table~\ref{tab:sub4kl}.

% The allocation study covers the 48 fused expert blocks, which hold 99.2\% of the quantized weight,
% with attention fixed at 4 bits and effective bits computed over the expert weight. The additive
% allocator prices each unit's 4-to-3 and 3-to-2 steps with 96 isolated damage measurements (each unit
% at 3 and at 2 bits, all others at 4) and greedily takes the step of minimum damage per bit saved;
% ModelOpt ranks units by squared-gradient sensitivity under the same greedy; the uniform mix demotes
% units in random order. The comparison is deliberately not measurement-matched: additive consumes the
% 96 measurements, ModelOpt only gradients, the uniform mix nothing. 

% The additive model ranks but does
% not price here: it predicts the KL of 40 random mixed configurations with $R^2 = 0.991$ (median
% relative error 7.8\%) yet underestimates the allocator's own chosen configurations by 16--57\%, an
% interaction the certificate forecast (at most 0.02 on these budgets) does not flag, so a quotable
% damage estimate at sub-4 budgets requires measuring the chosen configuration. 

In Table~\ref{tab:sub4kl}, the additive allocator's KL sits $3.0$--$6.6\times$ below ModelOpt's at
the servable budgets ($\ge 2.75$, established by the serving results below). CoopQ and the isolated
ranking, which reorders the additive allocator's own per-unit damages, occupy the same tier, while
MC-MoE and int8-outlier sit above ModelOpt. 
% The random uniform mix is not separated from the measured tier by KL alone, since it
% coincides with the all-3 configuration at eff 3.0 and has the lowest KL at 2.75, so the KL
% separation is from the gradient and prior scores, not from the random mix. 
% Below the servable range
% CoopQ leads, improving on the additive allocation by 8--28\%.

% A HAWQ-style ranking of the units by the same 96 isolated damages, without the per-bit pricing,
% produces allocations nearly identical to the additive allocator's (not shown), since the expert blocks
% have equal parameter counts, so damage ordering and damage-per-bit pricing
% nearly coincide. 

% Interaction-aware measurement pays only in the deep band: at effective bits 2.5 and
% 2.3, CoopQ and the allocators fitted on the 40 mixed configurations improve on the best 96-measurement
% allocation by 8--16\%, while at the servable budgets a 96-measurement allocation is best overall or
% within 1.3\% of best, at a fifth of CoopQ's measurement cost.
The construction replicates at scale on Qwen3-235B-A22B (Table~\ref{tab:sub4b235}): of the three
constructed allocations the additive one holds the lowest calibration KL at every sub-4 budget,
preserves reasoning when served at
effective bits 3.5, and degrades through 3.0.
 At effective
bits 2.9 the additive scoring rule appears on both sides of the survival boundary: derived on the
evaluation-task corpus itself, the additive and coverage allocations remain functional (42.9 and
41.2 AIME-2025), while the additive allocation of identical memory derived on a separately measured
corpus collapses (2.1, as does the random mix at 5.0). The gap opens before the boundary (58.3
against 39.6 at effective bits 3.0), so at this scale the derivation corpus, not the scoring rule,
sets the usable budget; every arm collapses at 2.75.

\input{generated/tab_sub4_235b}
% , with the additive
% configuration at effective bits 2.75 reaching calibration KL 0.310 against 1.008 for the uniform mix
% and 3.18 for ModelOpt.

\input{generated/tab_sub4kl}

% Serving these configurations tests whether the calibration-KL margins are realized as accuracy
% differences. We first me gsm8k (greedy decoding, 8{,}192-token generation cap, symbolic answer matching, 150
% problems per configuration, 300-problem bf16 anchor 96.7; the full KL-and-accuracy grid is
% Table~\ref{tab:sub4served}) they are, in three ways. At effective bits 3.5
% the additive configuration beats the ModelOpt configuration at identical memory, 97.3 against 91.3
% (all nine discordant problems favor it, McNemar $p = 0.0039$), the first budget in our study at which
% a matched-memory KL margin is realized as a served accuracy difference. At 2.9 and 2.75 the ModelOpt
% configurations stop terminating (the model stops emitting its end-of-sequence token, so every attempt
% runs to the token cap and yields no answer), while the additive configurations of the same memory
% solve the task (94.7 and 86.7), moving the smallest usable budget from about 3.0 effective bits under
% the gradient ranking to 2.75 under measured step damages. And the failure has a common boundary, a
% \emph{termination threshold}:  every configuration whose calibration KL exceeds roughly 0.45 nats stops
% terminating, however that KL was reached (the additive configuration at KL 0.448 also fails to
% terminate, while at 0.232 it solves the task), and degradation across this gsm8k grid is graded:
% generations terminate and solve the task below roughly 0.25 nats, produce coherent reasoning that
% fails to conclude within the cap near 0.45, and decay into repeated fragments beyond 0.6. Three
% configurations chosen
% before serving as an out-of-sample test of the threshold (the uniform mixes at effective bits 2.9 and
% 2.75, KL 0.18 and 0.25, and the all-3 configuration, 0.12) all terminate, at gsm8k 92.0, 88.7, and
% 95.3, so survival separates the measured-step allocation from the gradient ranking, not from the
% random mix.

% Table~\ref{tab:sub4served30} reports the full served grid. The methods that
% measure the ladder (additive, coverage with the measured ceiling, CoopQ, isolated-rank) occupy the top
% tier at every budget together with the uniform mix; the fitted surrogates lose 20--30 points at
% effective bits 3.0; and ModelOpt, MC-MoE, and int8-outlier collapse at 3.0 and below. Within the top
% tier the ordering moves with budget (additive leads at 3.5, coverage at 3.0, isolated-rank at 2.9, and
% the uniform mix at 2.75), so below four bits the decisive separation is tier membership, which tracks
% whether the method measured the ladder, rather than a single winner. 

Serving these configurations tests whether the calibration-KL margins are realized as served
accuracy. Table~\ref{tab:sub4code} reports HumanEval at both scales, and the full reasoning grid is
Table~\ref{tab:sub4served30} (App.~\ref{app:setup}). The failure mode has a common boundary, a
\emph{termination threshold}: above a model-specific calibration KL a configuration stops emitting
its end-of-sequence token, every generation runs to the token cap, and no answer is produced. On the
30B grid the boundary sits near 0.45 nats (\code{rf-alloc} at KL 0.449 collapses on AIME-2025 while
the additive configuration at 0.415 continues to score, Table~\ref{tab:sub4served30}), and the
threshold governs termination rather than accuracy. At effective bits 2.9 the four HumanEval methods
order exactly as their calibration KLs (0.215, 0.228, 0.236, and 1.042 giving 73.2, 44.5, 38.4, and
0.0). The accuracy gaps do not scale with the KL gaps:
CoopQ sits within 0.013 nats of the additive configuration yet 29 points below it, so which units
receive the 2-bit assignments
changes served accuracy at fixed memory and near-fixed KL. At 2.75 the ordering itself breaks: the
uniform mix has the lowest KL (0.319 against additive's 0.415) yet scores 19.5 against 54.9.


% gsm8k registers only the threshold (additive against uniform is $+5/-1$ discordants, not significant),
% so the accuracy contrast at matched survival needs a task with headroom below the threshold, and code
% provides it. On HumanEval (execution-graded, 164 problems, bf16 anchor 82.3) accuracy falls
% continuously along the additive configurations (68.9 at effective bits 2.9, 53.0 at 2.75), and at 2.9
% the allocators order exactly as their calibration KLs (additive 0.150, uniform 0.176, ModelOpt 0.445
% giving 68.9, 37.8, 0.0): the additive configuration is 31 points above the uniform mix at identical
% memory (56 of 61 discordant problems, McNemar $p < 10^{-10}$). Weight-only fake quantization assigns each unit a fixed
% tensor, so every configuration also exports as an ordinary bf16 checkpoint and serves natively; SGLang
% serving agrees with the Hugging Face harness on all nine configurations measured on both (identical
% orderings, accuracies within 4.3 points), 

% so Table~\ref{tab:sub4code} reports the full comparison on
% the production stack, with the Qwen3-235B replication. Two facts extend it. First, CoopQ separates
% calibration KL from served behavior: at effective bits 2.9 its configuration matches the additive one
% in KL (0.153 against 0.150) and in the multiset of levels (43 blocks at 3 bits, 5 at 2), yet scores
% 44.5 against 73.2, and at 2.75 it collapses to 6.1, below the uniform mix. Which units receive the
% 2-bit assignments therefore changes served accuracy at fixed memory and fixed calibration KL, and the
% threshold direction only survives one way: exceeding it collapsed every configuration measured, but
% CoopQ collapses at KL 0.226, well below the 0.45-nat boundary that held on the gsm8k grid, so a
% below-threshold KL does not guarantee termination across allocators. Second, the 235B serving was preregistered. The
% recorded predictions (the additive configuration at KL 0.310 terminates; the uniform mix at 1.008 and
% ModelOpt at 3.184 collapse; the three order as their KLs) held except termination: the additive
% configuration produces coherent code but runs past the cap on 101 of 164 problems, so the 0.45-nat
% threshold is not a constant across models. A second preregistration, recorded after this failure and
% before any further serving, placed the 235B onset in $(0.15, 0.31]$ and predicted the effective-bits
% 2.9 outcomes: the additive configuration (KL 0.150) terminates, the uniform mix (0.309) and ModelOpt
% (2.004) collapse. All three held (Table~\ref{tab:sub4code}), so the threshold transfers as an ordering
% with a model-specific value.

\begin{table}[t]
\centering\small
\caption{Served HumanEval accuracy (pass rate on base tests, all 164 problems, greedy decoding).
Cells at 0.0 are collapsed
configurations, with all or nearly all generations reaching the 8{,}192-token cap without terminating.
The two columns per model share one anchor measurement. The 235B CoopQ configurations
were not constructed due to measurement cost.}
\label{tab:sub4code}
\begin{tabular}{lcccc}
\toprule
 & 30B, eff 2.9 & 30B, eff 2.75 & 235B, eff 2.9 & 235B, eff 2.75 \\
\midrule
bf16 anchor & 83.5 & 83.5 & 88.4 & 88.4 \\
additive & \textbf{73.2} & \textbf{54.9} & \textbf{53.0} & \textbf{7.9} \\
CoopQ & 44.5 & 6.1 & --- & --- \\
uniform mix & 38.4 & 19.5 & 0.0 & 0.0 \\
ModelOpt & 0.0 & 0.0 & 0.0 & 0.0 \\
\bottomrule
\end{tabular}
\end{table}

% On the reasoning benchmarks under the model's recommended sampled decoding (AIME-2025, GPQA diamond,
% MMLU-Pro; Table~\ref{tab:sub4served30}), the same ordering reproduces at effective bits 2.9, and the
% additive and uniform configurations tie at 2.75.

% \subsection{Certificates, and where the margins appear}\label{sec:allocwhere}

% \paragraph{A certificate, not just a ranking.} A service-level constraint
% $\Delta\lce \le \varepsilon$ maps to a scalar budget on the summed log-headroom costs
% (App.~\ref{app:knapsack}), so the allocator returns a configuration together with a predicted
% $\Delta\lce$, quotable before committing ($\le 4\%$ error, Section~\ref{sec:interp}), which a ranking
% cannot produce. The ordinal certification of Section~\ref{sec:ordinal} supplies the matching upper
% bound: every allocator above scores configurations with some additive index, and the certified ordinal
% defect (the smallest fraction of pairwise ranking constraints any additive index must violate) is
% ${\sim}1$--$2\%$ in this regime, capping what any interaction-aware allocator can gain over additive
% scoring.

% \paragraph{Where the margins appear.} The matched-memory margin requires three conditions. The units
% must differ in sensitivity at the allocation granularity: the coefficient of variation of the
% in-context AS scores is $\approx 1$--$1.8$ on the MoE expert-layer inventories, with a
% few units carrying most of the sensitivity. The budget must be contended, low enough that some units
% must sit below their near-lossless precision. And the per-unit differences must survive the base
% quantizer, since a rotation that flattens them removes what allocation exploits
% (Section~\ref{sec:future}). The certificate plays a different role: by \thmone{}, it equals the best
% additive predictor's mean-squared error as a fraction of the damage variance, and it is not monotone in
% the margin (Qwen3-30B's $-17$ to $-21\%$ margin arises at a measured share of 0.013).

% \paragraph{A preregistered prospective test.} Before running any allocation on a fourth MoE
% (Qwen3.5-35B-A3B), we timestamped two predictions: that AS would attain lower KL than ModelOpt at every
% contended budget, and that the margin would fall in the $[-10,-40]\%$ band spanned by the three prior
% MoEs. The first was confirmed (Table~\ref{tab:35b}: AS best at every contended budget, with AS-additive
% and AS-coverage agreeing), and the band was rejected ($-5\%$ median), which identifies the covariate
% the band transfer was missing. This checkpoint is over an order of magnitude more KL-sensitive at
% matched eff (0.435 against Qwen3-30B's 0.032 at eff 4.5), and across the contended cells of the MoEs
% above, the AS-versus-ModelOpt margin shrinks as the shared damage floor grows: $-20/{-21}\%$ at eff
% 5.0/6.0 on Qwen3-30B ($\mathrm{KL}_{\mathrm{AS}} \approx 0.02$--$0.03$) against $-8.3/{-7.6}\%$ on the
% 35B ($\mathrm{KL}_{\mathrm{AS}} \approx 0.38$--$0.40$). A large shared damage floor, where every unit
% already hurts at its best available level, compresses what allocation can differentiate, so the margin
% band transfers only conditional on the damage level. The budget axis shows the same compression: at eff
% $\le 4.2$ every method makes the same few protection choices, freedom returns at 4.5--6.0 on this
% model, and the contended range itself moves with the damage floor.

% \begin{table}[t]
% \centering\small
% \caption{Qwen3.5-35B-A3B prospective test (KL, with the parenthetical giving the difference from
% AS-additive and $*$ marking paired significance over 16 shared sequences). Bold marks the lowest KL in
% the row. At eff 4.0 every greedy method degenerates to the identical all-4
% allocation, KL 0.631, and eff 4.2 sits close enough to that floor that no method separates.}
% \label{tab:35b}
% \begin{tabular}{ccccc}
% \toprule
% eff & naive-uniform & naive-routed & AS-add $=$ AS-cov & ModelOpt \\
% \midrule
% 4.2 & \textbf{0.458} ($-.005$) & 0.458 (off-budget: rule min eff $\approx 4.36$) & 0.462 & 0.465 ($+.003$) \\
% 4.5 & 0.457 ($+.022^*$) & 0.458 ($+.023$) & \textbf{0.435} & 0.449 ($+.014$) \\
% 5.0 & 0.430 ($+.030^*$) & 0.458 ($+.058^*$) & \textbf{0.400} & 0.436 ($+.037^*$) \\
% 6.0 & 0.378 ($+.002$) & 0.458 ($+.082^*$) & \textbf{0.375} & 0.406 ($+.030^*$) \\
% \bottomrule
% \end{tabular}
% \end{table}

% \begin{figure}[t]
%   \centering
%   \includegraphics[width=\linewidth]{fig7_alloc_pareto.png}
%   \caption{Allocation Pareto (fake-quant KL $\downarrow$ vs.\ effective bits) for Qwen3-30B (left) and
%   235B (right). AS-coverage and AS-additive attain the lowest KL of the compared methods at every 235B
%   budget and at the contended 30B budgets, and naive-uniform takes only the loose
%   235B eff-8 corner (near-lossless, nothing to allocate). Of the surrogates run as allocators,
%   \code{gp-alloc} trails AS by 3--18\% at the contended budgets, while \code{rf-alloc} is worse than
%   naive-uniform
%   (piecewise-constant marginals starve the greedy). ModelOpt falls below naive-uniform at eff 4.2/4.5 on
%   235B and at eff 8.0 on 30B.}
%   \label{fig:pareto}
% \end{figure}


\section{The reconciling theory: saturation generates both views}\label{sec:theory}

\subsection{The coverage model}\label{sec:coverage}

Section~\ref{sec:walsh} left two measurements in need of a common explanation: per-layer effects carry
most of the variance of the measured damage $f$, and some additive index reproduces at least $98\%$ of
its pairwise ranking of configurations. This section proposes saturation as the origin of both. We describe the damage with a saturating \emph{coverage model},
\begin{equation}
  f(S) = c\Bigl(1 - \prod_{i\in S}(1 - a_i)\Bigr),
  \label{eq:coverage}
\end{equation}
with a per-layer \emph{break-rate} $a_i\in[0,1]$ and a \emph{ceiling} $c>0$. Define the \emph{headroom}
$h(S) := c - f(S)$, the remaining distance to the ceiling. Imposing the coverage model is exactly the
statement that each quantized layer consumes a fixed fraction $a_i$ of whatever headroom remains,
$h(S\cup\{i\}) = (1-a_i)\,h(S)$ regardless of which layers came before.

This model is grounded in a measured property of the damage. Saturation, as we use the term, is the
conjunction of two facts visible on the configuration lattices of Section~\ref{sec:ordinal}: the loss
of a heavily quantized configuration is finite, and the marginal damage of one more quantized
layer shrinks as more layers are already quantized. Table~\ref{tab:marginals} reports the measurement
for each of the five certification cells, where a \emph{cell} is one model and quantizer setting
measured as a full $2^8$ lattice. Averaging $f(S\cup\{i\}) - f(S)$ over layers $i$
and contexts $S \not\ni i$ of a fixed size, we see that every cell's average marginal quantization
damage is largest when no other layer is quantized and falls once other layers are quantized
concurrently. The
evidence for coverage is assembled with each claim's standing in App.~\ref{app:evidence}.



Equivalently, the coverage model is the statement that $f$ is additive in log-headroom,
$-\log(1-f(S)/c) = \sum_{i\in S} b_i$ with $b_i:=-\log(1-a_i)$, hence coverage lies in the class
$g\bigl(\sum_{i\in S} b_i\bigr)$ of Section~\ref{sec:ordinal} with generator $g(u)=c(1-e^{-u})$. 
The measurements of Section~\ref{sec:ordinal} place $f$ within a ${\sim}2\%$
ordinal defect of this class across measured models. 
% What remains empirical about coverage is therefore only
% \emph{which} generator $g$ the data selects, not whether a $g\bigl(\sum b_i\bigr)$ form is legitimate.

% The names \emph{break-rate}, \emph{survival}, and \emph{noisy-OR} \citep{pearl1988probabilistic} come
% from the equivalent series-reliability reading of Eq.~\eqref{eq:coverage}, in which quantizing layer
% $i$ independently corrupts the prediction with probability $a_i$, a corrupted prediction incurs excess
% loss $c$, and $\prod_{i\in S}(1-a_i)$ is the probability the prediction survives. We use this reading
% for vocabulary and intuition, for instance the reciprocal-survival inflation factor of
% Section~\ref{sec:anchor}, while the claim the data validates is the headroom-composition rule above,
% not a per-token mixture. The evidence for coverage, spanning the mechanism above, the canonicity of the
% model class, the independent signed-expansion test, and the identification of the generator, is
% assembled with each claim's standing in App.~\ref{app:evidence}.

\subsection{The M\"obius and Walsh-Hadamard expansions}\label{sec:duality}

The coverage model was introduced to account for the variance and ranking measurements of
Section~\ref{sec:walsh}, and it is fitted to measured loss values (Section~\ref{sec:tauval}). It makes
a second prediction about a quantity that plays no role in that fit. The \emph{M\"obius/Harsanyi
transform} gives signed coordinates $\varphi(T)=\sum_{R\subseteq T}(-1)^{|T|-|R|}f(R)$ that
reconstruct the damage as $f(T) = \sum_{R \subseteq T} \varphi(R)$. 
Restricting the reconstruction sums to subsets of size $\le 1$
is precisely the isolate-and-sum aggregation of Section~\ref{sec:anchor}, since
$\varphi(\emptyset) = f(\emptyset) = 0$ and $\varphi(\{i\}) = f(\{i\})$. Restricting them to size
$\le 2$ adds exactly the pairwise finite differences
$\varphi(\{i,j\}) = f(\{i,j\}) - f(\{i\}) - f(\{j\})$ that CLADO \citep{clado2023} measures as its
cross-layer sensitivities.

\begin{customthm}{3}[M\"obius spectrum of coverage]\label{thm:mobius}
The M\"obius transform of the coverage model is, exactly,
$\varphi(T) = c\,(-1)^{|T|+1}\prod_{i\in T} a_i$ for $|T|\ge 1$ and $\varphi(\emptyset)=0$.
\end{customthm}

Coverage predicts a distinct form for this expansion. Its coefficients are nonzero at
every order and alternate in sign, so low-order truncations should overshoot and oscillate. 
The measured lattices show exactly this behavior
(Fig.~\ref{fig:mobius}). On Qwen3-0.6B at FP4 (8-layer exact lattice) the partial sums alternate, 
converging only at
$k=L$. On a 10-layer Qwen3-8B block with the full $2^{10}$-configuration lattice the same
divergence-then-cancellation appears at larger amplitude, swinging to $\pm 4.8$ against a true loss of
$0.148$, an overshoot of up to $32\times$. Because an order-$k$ M\"obius coefficient is an
alternating sum of up to $2^k$ measured values, the partial sums amplify measurement noise, so we
certify the swing by the split-half device of App.~\ref{app:ordinal}: partial sums computed from two
disjoint halves of the rounding draws independently reproduce the alternating overshoot (both swing
beyond $\pm 4$), while the same sums computed from the half-difference table, which contains noise
only, stay within $\pm 1.2$.
% This further shows that it is not sufficient to consider
% higher-order interaction terms in a predictor for the quantization damage.
% Taken alone, this behavior could be read as irreducibly
% high-order structure. Under coverage it is the expected consequence of saturation, produced by $L$
% per-layer parameters.

\begin{customthm}{4}[Walsh spectrum of coverage]\label{thm:walshspec}
Under the $p$-biased measure $\mup$, the Walsh coefficients are, exactly,
\[
  \fhat(T) = c\,(-1)^{|T|+1}\,(p(1-p))^{|T|/2}\Bigl(\prod_{i\in T} a_i\Bigr)\prod_{i\notin T}(1 - a_i p)
  \quad\text{for } T \ne \emptyset,
\]
and $\fhat(\emptyset) = \E_{\mup}[f] = c\bigl(1 - \prod_{i}(1 - a_i p)\bigr)$.
\end{customthm}

% The Walsh coefficient is the same $(-1)^{|T|+1}\prod a_i$ object as the M\"obius coefficient, multiplied
% by a
% geometric envelope and then squared in the energy $W_k$. That squaring is why the content that makes the
% signed series diverge appears in Section~\ref{sec:walsh} only as a small, cancellation-free $W_k$. The signed
% divergence and the variance concentration are one $L$-parameter model viewed signed versus squared. 
Proofs of both transforms are in App.~\ref{app:proofs}. 
% Each is a few lines, since the alternating subset-sum and
% the $\mup$-expectation of a product both factorize over coordinates, and both are verified to machine
% precision against brute-force enumeration.

\begin{figure}[t]
  \centering
  \includegraphics[width=0.9\linewidth]{fig1_mobius_divergence.png}
  \caption{The signed M\"obius expansion behaves as coverage predicts: order-$k$ truncations alternate and
  overshoot before full-order cancellation. Qwen3-8B 10-layer exact
  1024-configuration lattice, swings to $\pm 4.8$ (up to $32\times$ the true $0.148$), converging only at
  $k=10$ (right). The oscillation is certified against measurement noise: noise-only partial sums from the
  split-half difference table stay within $\pm 1.2$ (Section~\ref{sec:duality}).}
  \label{fig:mobius}
\end{figure}

\subsection{The saturation-depth law}\label{sec:tau}
By \thmone{}, the mean-squared error of the best additive predictor equals the order-$\ge 2$ energy
$\sum_{k\ge 2} W_k$. Dividing by the variance turns it into the order-$\ge 2$ share
\[
  \sharegeq := \frac{\sum_{k\ge 2} W_k}{\Var(f)},
\]
the fraction of the damage variance that per-layer effects leave unexplained, which is the normalized
form in which we report the certificate. 
The certificate
columns of the benchmark and allocation tables quote the share,
which is comparable across cells (Sections~\ref{sec:predictor} and \ref{sec:allocation}), while
statements in loss units use the absolute energies, namely the additive error floor of \thmone{} and
the sampling variance of the additive-coefficient estimator, which is the order-$\ge 2$ energy plus
noise (App.~\ref{app:setup}). 
On a
measured lattice the share is computed exactly, but at full-network scale, where no $2^L$ lattice can
be measured, it must be forecast from the $L+1$ fitted parameters of the coverage model. This
subsection derives that forecast, and that of the higher-order Walsh energies, from \propfive{} in
closed form. For a break-rate
profile $(a_1,\ldots,a_L)$ define the per-layer energy ratios and their aggregates
\[
  \beta_i := \frac{p(1-p)\,a_i^2}{(1-a_ip)^2},
  \qquad
  \tau^2 := \sum_i \beta_i,
  \qquad
  \Leff := \frac{\bigl(\sum_i\beta_i\bigr)^2}{\sum_i\beta_i^2} \in [1, L],
\]
with $\tau$ the \emph{saturation depth} and $\Leff$ the participation ratio of the profile, which
equals $L$ exactly when all $\beta_i$ are equal and approaches $1$ as one layer dominates. Let
$\es{k} := \sum_{|T|=k}\prod_{i\in T}\beta_i$ be the elementary symmetric polynomials in the
$\beta_i$.

\begin{customprop}{5}[heterogeneous certificate]\label{prop:hetcert}
Assume not all $a_i$ are zero. Under the coverage model,
\begin{enumerate}
\item[(i)] (product form) $W_k = c^2\bigl[\prod_i(1-a_ip)^2\bigr]\,\es{k}$, hence
  computable in closed form from the $L+1$ fitted parameters, which are themselves fit from
  $\mathcal{O}(L)$ measured configurations, and,
  \begin{equation}
    \sharegeq = 1 - \frac{\tau^2}{\prod_i(1+\beta_i) - 1};
    \label{eq:hetshare}
  \end{equation}
\item[(ii)] (leading order) $\sharegeq = \frac{\Leff-1}{2\Leff}\,\tau^2 + R$ with
  $|R| \le \tfrac14\tau^4 e^{\tau^2}$, so to second order in the per-layer amplitudes heterogeneity
  enters the certificate only through the effective depth $\Leff$;
\item[(iii)] at fixed $\tau^2$, $\sharegeq$ is Schur-concave in
  $(\beta_1,\ldots,\beta_L)$, hence
  \[
    \sharegeq \;\le\; 1 - \frac{\tau^2}{(1+\tau^2/L)^L - 1} \;\le\; 1 - \frac{\tau^2}{e^{\tau^2}-1},
  \]
  with equality on the left exactly when all $\beta_i$ are equal.
\end{enumerate}
\end{customprop}

\begin{proof}
%   \paragraph{Reading $\tau$.} Each Walsh coefficient is $|\fhat(T)| = c(1-ap)^L \beta^{|T|/2}$, so
% $\beta^{1/2} = \sqrt{p(1-p)}\,a/(1-ap)$ is the per-layer amplitude ratio: each additional participating
% layer scales the amplitude by $\beta^{1/2}$ and the energy by $\beta$. Normalizing the order-energies
% $\{W_k\}$ by $\Var(f)$ turns them into a probability distribution over the order $k$, describing where
% the variance lives. For the symmetric coverage model that distribution is exactly
% $\mathrm{Binomial}(L, q)$ with $q = \beta/(1+\beta)$, whose mean and variance
% both reduce to $L\beta$ for small $\beta$. Thus $\tau^2 = L\beta$ is, to leading order, the mean order of
% the spectrum and $\tau$ its standard deviation, linear in $a$ and growing as $L^{1/2}$ with depth, and the
% certificate is the tail statement $\sharegeq = P(K\ge 2 \mid K\ge 1)$ for $K$ so distributed, exact at
% finite $L$. The exact
% $\sharegeq(\tau)$ is the parameter-free curve of Fig.~\ref{fig:tau}, a sigmoid in $\tau$ with its
% transition at $\tau\approx 1$. The $\tfrac12\tau^2$ tail is its small-damage regime, and by \thmone{} it
% is also the best additive predictor's error floor. One consequence is worth stating here. For
% $\tau \lesssim 1$ the generator $g$ is traversed only over a nearly linear region, so the coverage and
% additive predictors nearly coincide, and the empirical near-tie between them is examined in
% Section~\ref{sec:predictor}.
For (i), \thmfour{} gives
$\fhat(T)^2 = c^2 (p(1-p))^{|T|} \prod_{i\in T} a_i^2 \prod_{i\notin T}(1-a_ip)^2$. Multiplying and
dividing by $\prod_{i\in T}(1-a_ip)^2$ factors every squared coefficient into the $T$-independent
constant $c^2\prod_{i=1}^{L}(1-a_ip)^2$ times $\prod_{i\in T}\beta_i$, so $W_k$ is that constant times
$\es{k}$, the variance $\sum_{k\ge 1} W_k$ is the constant times $\prod_i(1+\beta_i) - 1$ by the
generating identity $\prod_i(1+\beta_i) = \sum_{k\ge 0}\es{k}$. The claim follows.
Parts (ii) and (iii) are proven in App.~\ref{app:proofs}.
\end{proof}

%Part (i) is the form every certificate in this paper is evaluated from. 

% On blocks it is
% computed by exhaustive enumeration of \thmfour{} from the fitted $(c, a_i)$, and at full-network
% scale, where enumeration is impossible, the product form is what makes the certificate computable
% from fitted break-rates at all. 
% Part (iii) makes the symmetric case an exact
% worst case: a coverage-generated measurement can only sit on or below the envelope at matched
% $\tau$, so plotting measured shares against it (Fig.~\ref{fig:tau}) compares them with the exact upper
% bound. 
% At the equality case of part (iii), all break-rates equal, the certificate collapses to the
% one-parameter law $\sharegeq(\tau) = 1 - \tau^2/[(1+\tau^2/L)^L - 1]$, with small-$\tau$ limit
% $\frac{L-1}{2L}\,\tau^2 \approx \tfrac12\tau^2$ by part (ii) with $\Leff = L$. 
% This
% \emph{saturation-depth law} is the parameter-free curve of Fig.~\ref{fig:tau}.
% For a symmetric profile it is, by \thmone{}, the best additive
% predictor's error floor as a fraction of the damage variance, and for every other profile it bounds
% that fraction from above, by part (iii). 
Combining \thmone{} with part (ii), for small $\tau$ the
best additive predictor's mean-squared error is at most ${\approx}\tfrac12\tau^2\,\Var(f)$, so even
an exact coverage fit, whose error under the model is zero, can improve on additive prediction by no
more than that. 
% The benchmark and allocation results of Sections~\ref{sec:predictor}
% and~\ref{sec:allocation} exhibit the resulting near-tie once $N$ reaches a few multiples of $L$ (the
% two fits tie from $N \approx 4L$ in Fig.~\ref{fig:predn}), while at $N = L/2$ the additive fit's
% held-out error is roughly three times coverage's (0.315 against 0.093 on the 3B-wiki cell).
The symmetric case is stated as a corollary, with its binomial order
spectrum, in App.~\ref{app:proofs}.

% \paragraph{Reading $\tau$.} Each Walsh coefficient is $|\fhat(T)| = c(1-ap)^L \beta^{|T|/2}$, so
% $\beta^{1/2} = \sqrt{p(1-p)}\,a/(1-ap)$ is the per-layer amplitude ratio: each additional participating
% layer scales the amplitude by $\beta^{1/2}$ and the energy by $\beta$. Normalizing the order-energies
% $\{W_k\}$ by $\Var(f)$ turns them into a probability distribution over the order $k$, describing where
% the variance lives. For the symmetric coverage model that distribution is exactly
% $\mathrm{Binomial}(L, q)$ with $q = \beta/(1+\beta)$, whose mean and variance
% both reduce to $L\beta$ for small $\beta$. Thus $\tau^2 = L\beta$ is, to leading order, the mean order of
% the spectrum and $\tau$ its standard deviation, linear in $a$ and growing as $L^{1/2}$ with depth, and the
% certificate is the tail statement $\sharegeq = P(K\ge 2 \mid K\ge 1)$ for $K$ so distributed, exact at
% finite $L$.

% The bound runs in one direction only. A small $\tau$ caps the additive error for every break-rate
% profile, since parts (ii) and (iii) give $\sharegeq \lesssim \tfrac12\tau^2$ regardless of
% heterogeneity, but a large $\tau$ caps nothing, since by part (i) a profile concentrated on one
% layer has $\sharegeq = 0$ at any $\tau$. Certification therefore evaluates the product form of part
% (i) and uses $\tau$ only to certify smallness. 
% The same formulas locate the two regimes where the
% additive machinery fails. For a balanced profile ($\Leff \approx L$) at $\tau \gtrsim 1$, the
% one-parameter law puts the additive error above a third of the damage variance already at $\tau = 1$.
% For a concentrated profile ($\Leff \to 1$) the additive error is small, but all except one of the
% $L+1$ fitted coefficients then carry almost no variance, so at a fixed measurement budget their
% estimates are noise-dominated and per-layer prediction degrades (the 4B-wiki instance, participation
% ratio $\approx 1.2$ of 36 layers, Section~\ref{sec:scope}).

\subsection{Validation of the certificate}\label{sec:tauval}

\paragraph{The $\tau$-collapse (Fig.~\ref{fig:tau}).} We empirically test the saturation-depth law.
Each cell contributes one point, its measured order-$\ge 2$ share
against its fitted $\tau$, and \propfive(iii) requires any coverage-generated profile to fall on or
below the parameter-free
curve, with balanced break-rate profiles on it.
%  so that the single fitted scalar locates every cell
% between additive and many-body damage. 
We measured 36 full $2^8$ lattices, three models
(Qwen3-0.6B, Qwen2.5-3B, Llama-3.2-3B) crossed with twelve quantizer settings, integer
round-to-nearest at 8/6/5/4/3 bits with per-channel or per-tensor scales and 4-bit floating point
(FP4) with per-block or per-tensor scales (App.~\ref{app:setup}). For each cell we noise-correct the
order energies (App.~\ref{app:ordinal}), fit the coverage model to the same lattice by nonlinear
least squares (App.~\ref{app:setup}), and compare the measured order-$\ge 2$ share with the fit. The
measured share tracks the fitted $\tau = (\sum_i\beta_i)^{1/2}$ at
Spearman $\rho = 0.958$, and the exact share of \propfive(i), computed from the fitted $(c,a_i)$,
predicts the measured share at Pearson $\rho = 0.984$. Ten of the 36 cells exceed the curve, by at
most 8 percentage points. The largest exceedances are 8-bit cells whose absolute damage sits at the
measurement noise floor, where the noise-corrected share is unstable, and Llama-3.2-3B cells, the one
model with certified rank reversals (Table~\ref{tab:ordinal}), so the exceedances are consistent with
share noise plus the residual structure that coverage does not capture.
Repeating the comparison on Qwen3-30B-A3B
blocks (32 further cells, 8 block positions by 4 quantizer formats) gives Pearson $r = 0.997$.
% The large deviations in
% the figure run downward, the one direction \propfive(iii) permits. 
% The per-tensor cells fall
% 10--20\,pp below the prediction, consistent with true break-rate profiles more concentrated than the
% fit resolves (on the Qwen3-0.6B FP4 cell, order-1 terms carry 96.9\% of the measured variance but
% only 93.2\% of the fitted model's variance), and the few 4--8\,pp exceedances lie within the fit's
% 3.7\,pp median gap, in cells at the measurement noise floor or in 3-bit MoE cells below the
% fit-quality threshold. At small $\tau$ the lattice never approaches the ceiling, so only the
% products $c\,a_i$ are identifiable and the fitted $\tau$ there is noise-limited. One caveat applies.
% The grid uses a single calibration corpus, and the robustness of the fitted break-rates across
% corpora is established separately (Spearman $\rho = 0.74$--$0.78$ between break-rates fitted on wiki
% and on code, Section~\ref{sec:interp}).

\paragraph{Out-of-sample validation.} A certificate measured on a full lattice uses only \thmone{}.
When enumeration is intractable, we instead assume the coverage model and forecast the certificate
from break-rates fitted on sampled configurations. We
test the forecasting pipeline directly against held-out lattices, with per-cell results in
Table~\ref{tab:certest} of App.~\ref{app:setup}. The forecast classifies the near-additive and
heavily coupled regimes correctly (App.~\ref{app:setup}).

% We fit the break-rates on
% $N \approx L$ sampled configurations, the predictor's own budget, forecast the order-$\ge 2$ share
% with \propfive(i), and compare it with the share measured on the held-out full lattice. Near-additive
% cells forecast 2.6--6.4\% against measured 2.8--3.1\%, and many-body cells (interactions of every
% order among layers) forecast 59--63\% against 71--81\% (at most 4\,pp standard deviation over
% sampling seeds), so the forecast errs only by understating the share in the regime it already flags
% as beyond additive methods. On five exact lattices spanning true shares of 0.8\% to 92\%
% (Qwen3-0.6B and Qwen3-8B), the forecast also attains lower error at 14 of 15 cell-budget pairs than
% the
% two model-free estimators computable from the same configurations, the held-out residual of an
% additive fit and a split-half spectral estimate (Walsh coefficients from two disjoint halves of the
% configurations, multiplied so that independent noise cancels in expectation), with the per-cell
% detail in Table~\ref{tab:certest} of App.~\ref{app:setup}. 

\begin{figure}[t]
  \centering
  \includegraphics[width=\linewidth]{fig_tau_collapse.png}
  \caption{One parameter organizes the order-$\ge 2$ share. Left: measured (noise-corrected)
  order-$\ge 2$ energy share vs.\ fitted saturation depth $\tau = (\sum\beta_i)^{1/2}$ for 36 cells, 
  against the parameter-free symmetric law
  $\sharegeq(\tau)=1-\tau^2/[(1+\tau^2/L)^L-1]$. The curve is parameter-free: by
  \propfive(iii) it is the exact worst case at matched $\tau$ for a coverage-generated profile. Most
  per-channel and
  per-block cells sit on or below it (ten cells exceed it by at most 8 percentage points, noise-floor
  and Llama-3.2-3B cells, Section~\ref{sec:tauval}), while per-tensor 
  cells deviate below it, in the regime the theory flags as beyond coverage. Right: measured share vs.\ the exact
  heterogeneous \propfive{} prediction from the fitted $(c,a_i)$ (Pearson $r=0.984$; 30B-block extension
  $r=0.997$).}
  \label{fig:tau}
\end{figure}

\section{Scope}\label{sec:scope}

\paragraph{The activation channel carries the structure.} Measuring the damage set function on the
same layers with weights only (W4A16), activations only (W16A4), and both (W4A4) quantized, the weight
channel is non-saturating (isolate-and-sum error $\approx 0.10$, $\sum_i f(\{i\})/f(\text{all}) \approx 1.0$)
while the activation channel saturates strongly (isolate-and-sum 2.80 against coverage's 0.098 on the
3B W16A4 cell, Table~\ref{tab:seeds}). The saturation therefore lives in the activation channel,
which is why the method targets activation and W4A4 quantization.

\paragraph{Outlier-borne coupling under coarse activation scaling.} At the deployed NVFP4 per-block
format the order-$\ge 2$ share is 0.07 and the additive and coverage models attain the lowest errors,
while under per-tensor activation scales (a coarse setting that is not deployed) the share reaches
0.67 on the 0.6B certification lattice of Section~\ref{sec:ordinal} (the per-tensor cell of
Table~\ref{tab:ordinal}; the separately measured per-tensor lattice of Table~\ref{tab:certest}
puts its own share at 0.81), and the
nonparametric surrogates attain lower error, a breakdown the certificate forecasts. The coupling is
outlier-borne: localizing the scale (per-block-256)
drops the share to 0.01, clipping the top 0.1\% of activations drops it to 0.16, and SmoothQuant-style
scaling before quantizing drops it to 0.46. The fitted break-rates themselves are depth-structured
rather than outlier-driven (on Qwen3-0.6B they correlate with depth at Spearman $\rho = +0.68$ and
with output RMS at $\rho = +0.50$, but not with activation kurtosis). 

\paragraph{Measurement protocol.} Headline cells use 3 seeds (std $\le 0.03$), and the per-tensor
W4A4 isolate-and-sum failure is also 3-seed stable (0.6B $20.6\pm 0.3$, 3B $14.1\pm 0.3$), so the
rankings are robust.

% \paragraph{Cost is reported as memory.} Moving the always-active layers of the 30B MoE to FP8 changes
% decode latency by only $-2$ to $-5\%$ at concurrency 8--128, since MoE decode time is dominated by
% overhead and communication rather than by the GEMMs that quantization accelerates. We therefore report
% cost as effective bits, that is, as memory; realized speedup is a property of the serving stack.

\section{Related work}\label{sec:related}

\paragraph{Mixed-precision allocation via per-layer sensitivity.} AWQ \citep{lin2024awq} uses
activation-aware scaling, HAWQ and HAWQ-v2 \citep{hawq2019,dong2020hawqv2} a Hessian eigenvalue or trace,
BRECQ \citep{li2021brecq} and GPTQ \citep{frantar2022gptq} second-order weight reconstruction, and NVIDIA
ModelOpt's \code{auto\_quantize} \citep{modelopt}, the production method we benchmark, a Taylor/Fisher
score combined with a linear program over a sum of isolated per-layer scores. All of these aggregate
first- or second-order single-layer signals. In our deployment regime the closest neighbor is FGMP
\citep{fgmp2025}, which selects FP8-versus-NVFP4 blocks for weights and activations by a Fisher-weighted
perturbation score. Its score belongs to the same Taylor/Fisher family whose production embodiment,
ModelOpt, we benchmark throughout, while its intra-tensor block granularity requires co-designed kernels
that no serving stack we target provides, so it composes with layer-level allocation rather than
competing with it. ScaleBITS
\citep{scalebits2026} performs hardware-aligned bitwidth allocation for weights only.

\paragraph{Interaction-aware allocation.} CoopQ \citep{impq2025}, first circulated under the name IMPQ,
is the closest neighbor among interaction-aware methods. Its SPQE estimator samples layer permutations and progressively
quantizes along each, thereby measuring marginal contributions in context (permutation-sampled Shapley, at
a cost of roughly $100L$ quantized evaluations), then fits a linear-plus-pairwise surrogate
($a^\top q + q^\top K q$) and allocates by solving a binary quadratic program via MILP linearization.
CLADO \citep{clado2023} measures pairwise cross-layer effects by joint-quantization finite differences on
CNNs and ViTs and rejects the independence assumption.
Both fix the expansion order at two and do not test whether that order suffices. GuidedQuant
\citep{kim2025guidedquant} keeps within-output-channel dependencies but ignores cross-layer terms, and QEP
\citep{qep2025} models cross-layer accumulation as a compensation term. Sections
\ref{sec:walsh}--\ref{sec:theory} supply what none of these test, namely the order of the expansion, and
Section~\ref{sec:anchor} gives the closed-form account of the isolated-anchor bias they implicitly work around.

\paragraph{Configuration-quality surrogates.} For the prediction task itself, MiCo \citep{mico2025} fits
a random-forest accuracy predictor for edge-scale networks, AMQ \citep{amq2025} an RBF kernel model for
LLM weight-only quantization, and AutoQRA \citep{autoqra2026} a Gaussian-process surrogate inside a joint
search over bit-widths and adapter ranks. These predictor classes are accurate in-distribution but opaque,
and we show they degrade $2$--$5\times$ when extrapolating to more heavily quantized configurations
(Section~\ref{sec:predictor}).


% What
% that literature does not provide, and what quantization damage requires, is a saturating mechanistic
% generator together with the closed-form validity certificate it yields. Applied to quantization
% configurations the instrument is new in this domain, not in general.

\paragraph{}Additional related works are surveyed in App.~\ref{app:related}.

% Two adjacent strands are surveyed in App.~\ref{app:related}: MoE- and diffusion-specific mixed-precision
% allocators, against which the head-to-head comparisons appear in Section~\ref{sec:allocation}, and
% path-integral sensitivity methods, which we adapt to the layer-selection set function
% (Section~\ref{sec:allocation}, App.~\ref{app:as}). The literature on isolated-versus-in-context effects is
% discussed where it is used (Section~\ref{sec:anchor}).

% \paragraph{Position.} The surrogate task is crowded and the low-order-surrogate approach is established
% in adjacent domains. To our knowledge, the order-by-order measurement of quantization damage, the closed
% form connecting the divergent M\"obius series and the concentrated Walsh spectrum through one
% $L$-parameter model, and the self-certifying error theory are new, and the certificate is the piece none of
% the incumbent allocators or surrogates provides.

\section{Introduction}\label{sec:intro}

Quantizing a transformer to low bit-widths is the dominant route to cheap inference, and the formats now
being deployed quantize both weights and activations to 4 bits. What remains a design choice is
\textit{mixed precision}: which layers (or tensors) to keep at higher precision under a memory budget. The
choice is consequential, since allocations at the same memory budget can differ in output fidelity by tens
of percent (Section~\ref{sec:allocation}), and it is currently made by per-layer sensitivity heuristics.

These heuristics share a structural premise. Each scores a layer by a sensitivity measured largely in
isolation (i.e. with other layers unquantized), such as a Hessian eigenvalue or trace \citep[HAWQ;][]{hawq2019,dong2020hawqv2}, a
squared-gradient Fisher score \citep[ModelOpt;][]{modelopt}, or at most a pairwise interaction term
\citep{clado2023}, and then aggregates these scores to rank layers or to estimate the loss of a joint
configuration. To our knowledge, to what extent the loss of quantizing a \emph{set} of layers is in fact
recoverable from such low-order measurements has never been tested directly. Throughout, the \emph{order}
of a term is the number of layers it involves jointly, so per-layer terms are first-order and pairwise
terms second-order. The loss from quantization is a \emph{set function} defined on all $2^L$ subsets of $L$
layers (equivalently, the boolean cube $\{0,1\}^L$), which cannot be enumerated for models of any realistic depth.

To measure the order of this function directly, we write
$f(S) = \lce(S\ \text{quantized}) - \lce(\text{full precision})$ for the damage of quantizing the layer set
$S$, obtained by evaluating the partially quantized model, and analyze the measured $f$ through two
classical changes of basis. The Walsh transform (the functional-ANOVA decomposition) splits the variance
of $f$ by order under the \emph{deployment distribution}, the distribution in which each layer is
quantized independently with probability $p$. The M\"obius expansion reconstructs $f$ from signed terms of
increasing order. This analysis yields two findings.

The first finding is that the variance of the damage is almost entirely first-order. Across configurations
drawn from the deployment distribution, 85--93\% of the variance of $f$ is explained by per-layer effects,
rising to 98--99\% on exact lattices (blocks of $L$ layers for which every one of the $2^L$ configurations
is measured) once measurement noise is corrected for (Fig.~\ref{fig:walsh}, Section~\ref{sec:walsh}). 
The unexplained remainder therefore serves as a \emph{certificate}: a single number, measured on full
lattices (an identity, by \thmone{}) or forecast in closed form from fitted break-rates at
full-network scale
(Section~\ref{sec:tau}), that states in advance
how well any additive model can do, and which we report alongside
every result. The additive model attaining this error floor also has a closed form
(Section~\ref{sec:anchor}).
% A simple
% diagnostic then locates what fails in the incumbent heuristics. Aggregating isolated sensitivities is
% equivalent to fitting an additive model constrained to predict zero damage at the unquantized network. On
% the exact lattice where FP4 damage is large, that constrained fit reaches $R^2 = 0.13$, and freeing the
% single intercept parameter raises it to $0.94$ (Section~\ref{sec:walsh}). The additive model class therefore
% contains accurate predictors, and the failure lies in the zero-damage constraint that isolated scores
% impose.

The second finding is that the variance of $f$ explained by higher-order terms does not reflect a failure of per-layer structure. The
order-$\ge 2$ share, the fraction of the variance that per-layer effects leave unexplained, equals by
\thmone{} the mean-squared error of the best additive fit divided by the variance
(Section~\ref{sec:walsh}). This share can be nonzero
even when the damage is perfectly described by per-layer terms. If $f$ is additive after a strictly increasing
transformation, that is, it has the form $g\bigl(\sum_{i\in S} b_i\bigr)$ for per-layer contributions
$b_i$ and a strictly increasing $g$, then composing with $g^{-1}$ yields a function whose
order-$\ge 2$ energy is exactly zero. Whether $f$ belongs to this class is testable on rankings alone
(\thmtwo), and we measure the smallest fraction of pairwise configuration comparisons that any additive
index must order incorrectly, certified against measurement noise (Section~\ref{sec:ordinal}). This
fraction is at most 2\%, on some models indistinguishable from zero, and a direct search over
transformations agrees. The
measured damage is therefore of the form $f(S) \approx g\bigl(\sum_{i\in S} b_i\bigr)$.
This model class recurs throughout the paper. It justifies the form fitted in Section~\ref{sec:theory}, and it
caps how much any interaction-aware allocator can improve on additive configuration scoring
(Section~\ref{sec:allocation}).

We show that a single property, \emph{saturation}, explains both findings exactly. Because the loss is
bounded, one can define the \emph{headroom} $h(S) = c - f(S)$, the damage still available below a ceiling
$c$. The measurements are well described by a model in which each quantized layer consumes a fixed fraction
$a_i$ of the remaining headroom, $h(S\cup\{i\}) = (1-a_i)\,h(S)$, independently of which layers preceded
it. Equivalently, $f(S)=c\bigl(1-\prod_{i\in S}(1-a_i)\bigr)$, which we call the \emph{coverage model}.
The model is fitted to the measured losses of sampled
configurations by least squares. 

% Expanding the fitted model in the M\"obius basis then predicts a property
% of coefficients that played no role in the fit: they should be nonzero at every order and alternate in
% sign, so low-order truncations of the measured expansion should overshoot and oscillate.
% The exact lattices confirm this behavior, with partial sums overshooting the true loss by up to
% $32\times$ before cancelling at full order (Fig.~\ref{fig:mobius}, Section~\ref{sec:duality}). 
% Taken alone, such
% divergence would suggest that no truncation of the expansion at any low order can describe $f$. Under
% coverage it is the expected consequence of saturation, produced by $L$ per-layer parameters. 
% Nor is 
The model class of the second finding is not arbitrary. By a representation theorem of Acz\'el and Ling
\citep{aczel1966fe,ling1965assoc}, any rule that accumulates damage one layer at a time through a fixed
operation that is associative, continuous, and strictly increasing can be rewritten in the form
$g\bigl(\sum_i b_i\bigr)$ (App.~\ref{app:link}). 
% Assuming such a model therefore assumes only that damage accumulates by a
% fixed rule of this kind, and coverage is the member of this class that the measured headroom
% consumption selects (App.~\ref{app:link}).

% This structure supports two predictors of a configuration's loss, each with $L+1$ parameters, and an
% accuracy statement for them that is exact rather than an upper bound. 
% The \textbf{additive} model is the
% optimal first-order predictor, and by Parseval's identity
% its mean-squared error equals the order-$\ge 2$ Walsh energy of $f$, the portion of the variance carried by
% terms above first order (\thmone). That
% quantity therefore serves as a \textbf{certificate}: a single number, measured on full lattices or computed
% in closed form from fitted break-rates at full-network scale (Section~\ref{sec:tau}), that states in advance
% exactly how well any additive model can do, and which we report alongside every result. 

% Both predictors, the fitted coverage model and the additive model, predict a held-out
% configuration's loss to within roughly 6--13\% from about $L$ measurements, and the two predict comparably,
% which the theory explains: no predictor can improve on the best additive one by more than the
% order-$\ge 2$ energy, and at deployed damage levels that energy is a few percent of the variance
% (Section~\ref{sec:tau}). The certificate also identifies when
% neither suffices, as an order-$\ge 2$ share near $0.3$ marks the point at which nonparametric surrogates
% become preferable (Section~\ref{sec:predictor}).

Aggregating isolated sensitivities is equivalent to fitting an additive model constrained to predict
zero damage at the unquantized network, and this constraint is what fails in the incumbent
heuristics: freeing the single intercept parameter raises $R^2$ from $0.13$ to $0.94$ on the exact
lattice where FP4 damage is large (Section~\ref{sec:walsh}). Because saturation shrinks a layer's
marginal damage as more layers are quantized, a sensitivity measured with every other layer at full
precision is inflated by a closed-form factor, the quantitative form of a context dependence known
qualitatively
\citep{ancona2020shapley,mcgrath2023hydra,impq2025}, and summing $L$ inflated coefficients compounds
the error to 71--376\%, which we measure for isolate-and-sum (i.e. the quantization loss is the sum of the damages from quantizing each layer separately)
 and for a Hessian-trace HAWQ-v2
(Section~\ref{sec:anchor}). A single
\emph{Aumann--Shapley path integral} recovers the additive model's optimal coefficient evaluated at a
layer-specific effective deployment density, and with it the exact ranking of layers, with no sampled
configurations, a recovery we prove under the coverage model (Section~\ref{sec:predictor},
App.~\ref{app:as}).

Because both models predict rather than rank, they can \emph{price} a configuration, that is, quote its
loss before committing to it, which a ranking cannot do. On unseen configurations that protect the most
fragile layers and quantize the rest, the quoted loss is
accurate to within 4\% (App.~\ref{sec:interp}). Against the black-box surrogates used for the same task,
random forests \citep{mico2025}, Gaussian processes \citep{autoqra2026}, and RBF kernels \citep{amq2025},
the $(L{+}1)$-parameter models are best or tied at small measurement budgets, the regime that matters when
every measurement is a forward pass. The strongest of these, the Gaussian process, succeeds for a reason
our analysis identifies and fails where the theory predicts, degrading $2$--$5\times$ under
extrapolation to differently quantized configurations than those it was fit on, where the additive and coverage
models stay within $2\times$. That direction of failure matters because
an allocator must price configurations quantized at various levels
(Section~\ref{sec:allocation}).

Used as allocators at matched memory, the same fitted models attain the lowest KL divergence among the
compared allocators (sensitivity scores, structural rules, and the black-box surrogates converted into
allocators) on mixture-of-experts language models across models from 0.6B to 355B parameters
(Section~\ref{sec:allocation}). Below four bits configurations allocated from the fitted models
continue to solve code and reasoning tasks at effective bit-widths where configurations
from other methods of the same memory budget no longer produce terminating generations, at both 30B and 235B
(Section~\ref{sec:allocsub4}).

\paragraph{Contributions.}
\begin{enumerate}
% \item An exact \textbf{Walsh (functional-ANOVA) variance decomposition} of low-bit quantization damage
%   on full $2^L$ lattices: 85--93\% of the variance is first-order, and the coordinate-free ordinal defect
%   is at most 2\% (Section~\ref{sec:walsh}, Fig.~\ref{fig:walsh}).
\item A \textit{closed-form theory}: the coverage model's exact M\"obius and Walsh spectra
  (Theorems~3 and~4), corroborated by the confirmed divergence of the signed expansion
  (Fig.~\ref{fig:mobius}); the certificate identity (\thmone) with an exact $O(L)$ heterogeneous form
  (\propfive); and the ordinal characterization of additivity up to an increasing transformation
  (\thmtwo)
  (Sections~\ref{sec:walsh}--\ref{sec:theory}, App.~\ref{app:proofs}).
\item A \textit{negative result with a closed-form account}: sensitivities measured in isolation,
  including HAWQ-v2, mis-price a configuration's loss by 71--376\%, and the inflation factor
  responsible is derived in closed form (Section~\ref{sec:anchor}).
\item A \textit{fit-free estimator}: one Aumann--Shapley path integral recovers the additive model's
  optimal coefficient at a layer-specific effective density, and its layer ranking exactly, proven
  under the coverage model (Section~\ref{sec:predictor}, App.~\ref{app:as}).
\item An \textit{interpretable predictor} that is best or tied against RF/GP/RBF/pairwise baselines at
  small budgets, degrades at most $2\times$ under configuration extrapolation where the GP
  degrades
  $2$--$5\times$, and is validated on LLMs 0.6B--355B and a diffusion DiT
  (Section~\ref{sec:predictor}, Fig.~\ref{fig:ood}).
\item An \textit{allocator at matched memory} that attains the lowest KL among sensitivity-, rule-,
  and surrogate-based allocators (17--21\% below the production ModelOpt baseline on Qwen3-30B,
  13--27\% on Qwen3-235B, 20--37\% on GLM-4.6) and decides served task survival below four bits.
\end{enumerate}

% \paragraph{Scope.} The order-energy measurements and the predictor concern activation quantization at
% FP4 / W4A4, the deployed regime. Weight-only quantization is near-additive, and KV-cache quantization more
% so, so neither carries the structure studied here. The prediction results span LLMs from 0.6B to 355B
% parameters and two diffusion DiTs, while the allocation results hold at matched memory in the
% contended-budget regime, with their limits mapped in Section~\ref{sec:scope}.
